\section{Sharp local operator algebra and Lane~A reconstruction}
\label{sec:sharp_local}

This section is the constructive Lane~A packet on the frozen proof spine.
Its output is used twice in the final claim: Appendix~\ref{app:axioms_package} consumes the inverse-flow, quotient-block, operator-domain, and cyclicity statements proved here in order to package the OS/Wightman existence and non-triviality theorem; Appendix~\ref{app:global_form_resolution} then uses the same Lane~A output to identify the exact local field content of the theorem-level endpoint.
The OS/Wightman-strength axioms themselves are therefore promoted to a separate claim-facing packet in Appendix~\ref{app:axioms_package}; the present section supplies the constructive statements that make that packet load-bearing.

\subsection{Auxiliary: quasi-locality of flowed probes}
This section records the \emph{flow-time quasi-locality and derivative bounds}
for the flowed composites $\mathcal O_{S,t}$ and $\mathcal O_{Q,t}$, in a form uniform in $a\to 0$ on the AF window.
The deterministic hypotheses used here (a local chart/boundedness regime for the DeTurck flow) are verified on the
good-flow chart event in Section~\ref{sec:good_flow_chart} (Lemma~\ref{lem:chart_persist}), and the contribution of the
complement is absorbed in expectations by truncation stability (Corollary~\ref{cor:superpoly_flow}).

\paragraph{Important correction (hardening).}
The gradient flow is \emph{parabolic}; it does \emph{not} have finite propagation speed.
Thus ``strict support in $B(x,c\sqrt t)$'' is false. What is true (and sufficient for the dyadic polymer/RG framework) is:
\begin{itemize}
\item \textbf{quasi-locality}: influence of initial data at distance $r$ is suppressed like $\exp(-c r^2/t)$;
\item \textbf{derivative bounds}: functional derivatives of the flow map and of $\mathcal O_{S,t},\mathcal O_{Q,t}$ satisfy Gaussian off-diagonal bounds with
explicit $t$-scaling factors;
\item \textbf{local truncation}: by cutting off outside a ball of radius $R\sqrt t$ one incurs an error $\lesssim e^{-cR^2}$, which can be made $\le t^\eta$ by choosing $R\sim \sqrt{\log(1/t)}$.
\end{itemize}

This module replaces 's strict-support formulation with the correct quasi-local statement and makes the constants explicit.

We work in Euclidean $\mathbb R^4$ (or a periodic box of side $L$ with $t\le t_\ast\ll L^2$ so boundary effects are negligible).
Let $A_\mu(x)\in\mathfrak g$ be the initial gauge potential and $B_\mu(t,x)$ the gradient-flowed field for $t>0$.

The standard (continuum) Yang--Mills gradient flow is
\begin{equation}\label{eq:cov_flow}
\partial_t B_\mu \ =\ D_\nu G_{\nu\mu},\qquad B_\mu(0,\cdot)=A_\mu,
\end{equation}
where $D_\nu=\partial_\nu+[B_\nu,\cdot]$ and $G_{\mu\nu}=\partial_\mu B_\nu-\partial_\nu B_\mu+[B_\mu,B_\nu]$.

To access parabolic estimates we use the DeTurck gauge-fixing term:
\begin{equation}\label{eq:deturck_flow}
\partial_t B_\mu \ =\ D_\nu G_{\nu\mu} + D_\mu(\partial_\nu B_\nu),\qquad B_\mu(0,\cdot)=A_\mu.
\end{equation}

\begin{proposition}[Gauge equivalence and invariance of gauge-invariant composites]\label{prop:deturck_equiv}
Solutions of \eqref{eq:cov_flow} and \eqref{eq:deturck_flow} are related by a (time-dependent) gauge transformation $g(t,x)\in G$:
\[
B^{\mathrm{cov}}_\mu(t)=g(t)\,B^{\mathrm{det}}_\mu(t)\,g(t)^{-1}-\partial_\mu g(t)\,g(t)^{-1}.
\]
Consequently, gauge-invariant local observables built from $G_{\mu\nu}(t)$, in particular
$\mathcal O_{S,t}(x)=\tfrac14\mathrm{tr}\,G_{\mu\nu}(t,x)G_{\mu\nu}(t,x)$ and
$\mathcal O_{Q,t}(x)=\tfrac14\mathrm{tr}\,G_{\mu\nu}(t,x)\widetilde G_{\mu\nu}(t,x)$,
have identical values under the two flows.
\end{proposition}

\paragraph{Comment.}
This is the standard DeTurck trick for Ricci/YM flows: the gauge-fixing term generates a compensating gauge transformation.
Thus we may work with \eqref{eq:deturck_flow} without loss for gauge-invariant observables.

To obtain uniform constants we work on a chart-controlled regime (consistent with the patched RG philosophy of Section~\ref{sec:patching}).

\begin{proposition}[Bounded-regime hypothesis (verified on the good-flow chart event)]\label{ass:bounded_regime}
Fix $t_\ast>0$. On the good-flow chart event $G_{\Lambda}(\varepsilon)$ of Definition~\ref{def:good_plaq} with initial chart extraction as in
Lemma~\ref{lem:tree_chart_local}, the DeTurck-flow solution satisfies the scale-consistent bound
\begin{equation}\label{eq:bounds_B}
\sup_{0<t\le t_\ast}\Big(\|B(t)\|_{L^\infty}+ t^{1/2}\|\nabla B(t)\|_{L^\infty}\Big)\ \le\ M_0+M_1.
\end{equation}
In lattice language, this corresponds to a fixed exponential chart bound on link angles and plaquette angles (e.g.\ $\le\delta<\pi/4$) on the relevant scales.
\end{proposition}

\paragraph{Where this comes from.}
Lemma~\ref{lem:chart_persist} (Section~\ref{sec:good_flow_chart}) gives a deterministic sufficient condition for
\eqref{eq:bounds_B}: on $G_{\Lambda}(\varepsilon)$ the flow stays inside a fixed exponential chart up to time $t_\ast$ and the parabolic estimates
control $\|B(t)\|_\infty$ and $t^{1/2}\|\nabla B(t)\|_\infty$ uniformly on $(0,t_\ast]$.
In the main proof we invoke reflection positivity/OS only on unflowed positive-time observables; when flowed insertions appear in expectations,
we enforce $G_{\Lambda}(\varepsilon)$ locally and absorb the complement by Corollary~\ref{cor:superpoly_flow}.

Let $\mathcal F_t$ denote the DeTurck flow map $A\mapsto B(t)$.
Let $\psi_\mu(t,x)$ denote the first variation of $B_\mu(t,x)$ under an infinitesimal perturbation of the initial data $A$.

\begin{lemma}[Linearized DeTurck flow is uniformly parabolic]\label{lem:lin_parabolic}
Under Proposition~\ref{ass:bounded_regime}, the first variation $\psi$ satisfies a linear parabolic system of the form
\begin{equation}\label{eq:lin_system}
\partial_t \psi \ =\ \Delta \psi + \sum_{\alpha} a_\alpha(t,x)\,\partial_\alpha \psi + b(t,x)\psi,
\end{equation}
where $\Delta$ is the Euclidean Laplacian acting componentwise and the coefficients satisfy the uniform bounds
\[
\begin{aligned}
\|a_\alpha(t,\cdot)\|_{L^\infty(\mathbb R^4)}&\le C_{a,0}(M_0),\\
\|b(t,\cdot)\|_{L^\infty(\mathbb R^4)}&\le C_{b,0}(M_0)+C_{b,1}(M_0,M_1)\,t^{-1/2},
\qquad 0<t\le t_\ast.
\end{aligned}
\]
\end{lemma}

\paragraph{Meaning.}
Equation \eqref{eq:lin_system} is uniformly parabolic; on any time slab $[s,t]\subset(0,t_\ast]$ the coefficients are bounded.
The possible $t^{-1/2}$ singularity at $t\downarrow 0$ is time-integrable and is treated by the parametrix/Volterra majorant
in Appendix~\ref{app:kernel_parametrix} (and uniformly in $a$ in Appendix~\ref{app:discrete_parametrix}).

\begin{proposition}[Gaussian kernel bounds for the linearized operator (proved in Appendix~\ref{app:kernel_parametrix} and Appendix~\ref{app:discrete_parametrix})]\label{ass:gaussian_kernel}
Let $K(t,s;x,y)$ be the fundamental solution (matrix kernel) of \eqref{eq:lin_system} on $0\le s<t\le t_\ast$.
Assume there exist constants $c_0,c_1>0$ and for each integer $m\ge 0$ constants $C_m$ (depending only on $d=4$ and the coefficient bounds
$C_{a,0},C_{b,0},C_{b,1}$ (and the corresponding finite set of derivative bounds used in the parametrix)) such that:
\begin{align}
\|K(t,s;x,y)\|_{\mathrm{op}}
&\le c_0\,(t-s)^{-2}\exp\!\Big(-c_1\,\frac{|x-y|^2}{t-s}\Big),\label{eq:K_gauss}\\
\|\nabla_x^m K(t,s;x,y)\|_{\mathrm{op}}
&\le C_m\,(t-s)^{-(2+m/2)}\exp\!\Big(-c_1\,\frac{|x-y|^2}{t-s}\Big).\label{eq:K_der_gauss}
\end{align}
\end{proposition}

On the lattice, an analogous statement holds for the discrete operator under uniform chart bounds; constants can be chosen independent of $a$ for $t\ge c a^2$.

Let $\delta A$ be a perturbation supported near a point $y$.

\begin{theorem}[First-variation quasi-locality of the flow map]\label{thm:first_var}
Under the bounded-regime hypothesis~\ref{ass:bounded_regime} and the Gaussian kernel bounds~\ref{ass:gaussian_kernel},
for $t\in(0,t_\ast]$ the Fr\'echet derivative of the flow map satisfies the kernel bound:
\[
\left\|\frac{\delta B_\mu(t,x)}{\delta A_\nu(y)}\right\|
\ \le\ c_0\,t^{-2}\exp\!\Big(-c_1\,\frac{|x-y|^2}{t}\Big),
\]
and for spatial derivatives:
\[
\left\|\nabla_x^m\frac{\delta B_\mu(t,x)}{\delta A_\nu(y)}\right\|
\ \le\ C_m\,t^{-(2+m/2)}\exp\!\Big(-c_1\,\frac{|x-y|^2}{t}\Big).
\]
Constants are those of Proposition~\ref{ass:gaussian_kernel} with $s=0$.
\end{theorem}

\begin{corollary}[First-variation bounds for field strength]\label{cor:first_var_G}
Under the same assumptions, the field strength variation satisfies for $m\ge 0$:
\[
\left\|\nabla_x^m\frac{\delta G_{\rho\sigma}(t,x)}{\delta A_\nu(y)}\right\|
\ \le\ C'_m\,t^{-(5/2+m/2)}\exp\!\Big(-c_1\,\frac{|x-y|^2}{t}\Big),
\]
where $C'_m$ depends on $C_{m+1}$ and $M_0$ (through commutator terms in $G$).
\end{corollary}

The polymer seminorms used in dyadic RG control derivatives up to a fixed finite order $r_{\max}$.
We therefore state a finite-order induction hypothesis.

\begin{proposition}[Finite-order higher derivative bounds (proved in Appendix~\ref{app:kernel_parametrix})]\label{ass:higher_der}
Fix $r_{\max}\ge 1$.
Assume that for each $1\le r\le r_{\max}$ there exist constants $C^{(B)}_{r,m}$ such that the $r$-th Fr\'echet derivative of the flow map obeys:
\begin{equation}\label{eq:high_der_B}
\left\|\nabla_x^m\frac{\delta^r B(t,x)}{\delta A(y_1)\cdots\delta A(y_r)}\right\|
\ \le\ C^{(B)}_{r,m}\, t^{-(2+m/2+r)}\,
\exp\!\Big(-c_1\,\frac{\mathrm{dist}(x;\{y_i\})^2}{t}\Big),
\end{equation}
where $\mathrm{dist}(x;\{y_i\})=\min_i |x-y_i|$.
\end{proposition}

\paragraph{Comment.}
The exponent $2+m/2+r$ is a safe (nonoptimal) power count; any explicit polynomial bound in $t^{-1/2}$ is sufficient for the A1 insertion norms,
because $r_{\max}$ is fixed and the ultimate SFTE remainder is controlled by extracting relevant/marginal parts and contracting irrelevants.
Proposition~\ref{ass:higher_der} is proved by differentiating the flow equation $r$ times and applying Duhamel/Gr\"onwall with the Gaussian kernel bounds.

\begin{corollary}[Higher derivative bounds for $\mathcal O_{S,t}$ and $\mathcal O_{Q,t}$]\label{cor:high_der_OSOQ}
Assume~\ref{ass:bounded_regime},~\ref{ass:gaussian_kernel}, and~\ref{ass:higher_der}.
Then for each $r\le r_{\max}$ there exist constants $C^{(S)}_{r}$ and $C^{(Q)}_{r}$ such that
\begin{align}
\left\|\frac{\delta^r \mathcal O_{S,t}(x)}{\delta A(y_1)\cdots\delta A(y_r)}\right\|
&\le C^{(S)}_{r}\, t^{-(5+r)}\,
\exp\!\Big(-c_1\,\frac{\mathrm{dist}(x;\{y_i\})^2}{t}\Big),\label{eq:der_OS}\\
\left\|\frac{\delta^r \mathcal O_{Q,t}(x)}{\delta A(y_1)\cdots\delta A(y_r)}\right\|
&\le C^{(Q)}_{r}\, t^{-(5+r)}\,
\exp\!\Big(-c_1\,\frac{\mathrm{dist}(x;\{y_i\})^2}{t}\Big).\label{eq:der_OQ}
\end{align}
The constants depend on $M_0,M_1$ through the coefficient bounds and on the fixed derivative order $r$.
\end{corollary}

\paragraph{Remark.}
The power $5+r$ is again safe (nonoptimal). The key structural feature is the Gaussian suppression $\exp(-c\,\mathrm{dist}^2/t)$.

We now convert quasi-locality into a localization statement usable in the dyadic polymer framework.

\begin{definition}[Ball truncation radius]\label{def:Rt}
For $R\ge 1$ define the ball $B_R(x,t):=\{y:|y-x|\le R\sqrt t\}$.
For a perturbation $\delta A$ supported in $B_R(x,t)^c$ define the Lipschitz seminorm
\[
\mathrm{Lip}^{(R)}_x(\mathcal O_{*,t})
:=
\sup_{\substack{\delta A\neq 0\\ \mathrm{supp}(\delta A)\subset B_R(x,t)^c}}
\frac{|\mathcal O_{*,t}(x;A+\delta A)-\mathcal O_{*,t}(x;A)|}{\|\delta A\|_{L^\infty}}.
\]
\end{definition}

\begin{theorem}[Exponential tail outside $B_R(x,t)$]\label{thm:tail_R}
Under the first-variation bound of Theorem~\ref{thm:first_var} and the corresponding bound for $G$ (Corollary~\ref{cor:first_var_G}),
the following exponential tail estimate holds.
There exist constants $C_{\mathrm{tail}},c_{\mathrm{tail}}>0$ such that for $*\in\{S,Q\}$,
\[
\mathrm{Lip}^{(R)}_x(\mathcal O_{*,t})\ \le\ C_{\mathrm{tail}}\,t^{-5}\,\exp(-c_{\mathrm{tail}}R^2).
\]
Consequently, if two initial data sets agree on $B_R(x,t)$, then
\[
|\mathcal O_{*,t}(x;A)-\mathcal O_{*,t}(x;A')|
\le C_{\mathrm{tail}}\,t^{-5}\,\exp(-c_{\mathrm{tail}}R^2)\,\|A-A'\|_{L^\infty}.
\]
\end{theorem}

\begin{proof}
Integrate the first functional derivative along the straight line path $A_\theta=A+\theta(A'-A)$ and use the Gaussian tail:
\[
\mathcal O_{*,t}(A')-\mathcal O_{*,t}(A)=\int_0^1 \int \frac{\delta \mathcal O_{*,t}}{\delta A(y)}(A_\theta)\,(A'-A)(y)\,dy\,d\theta.
\]
Restrict $y\in B_R^c$ and bound the derivative kernel by \eqref{eq:der_OS} or its $r=1$ version, then integrate the Gaussian tail:
$\int_{|x-y|>R\sqrt t} t^{-2}e^{-c|x-y|^2/t}dy\lesssim e^{-c'R^2}$.
\end{proof}

\begin{corollary}[Choosing $R(t)$ to match a target power $t^\eta$]\label{cor:R_choice}
Fix $\eta>0$. Choose
\[
R(t):=\sqrt{\frac{(\eta+6)\log(1/t)}{c_{\mathrm{tail}}}}\qquad (t\in(0,e^{-1}]).
\]
Then $t^{-5}e^{-c_{\mathrm{tail}}R(t)^2}\le t^{\eta+1}$.
Hence the truncation error outside $B_{R(t)}(x,t)$ is $\le C_{\mathrm{tail}} t^{\eta+1}\|A-A'\|_\infty$.
\end{corollary}

\paragraph{Interpretation for A1.}
Corollary~\ref{cor:R_choice} shows that, for any fixed target exponent $\eta$ (e.g.\ $\eta=1$ suffices for the SFTE remainder bookkeeping in Theorem~\ref{thm:ins_decomp}), one may treat $\mathcal O_{S,t}(x)$ and $\mathcal O_{Q,t}(x)$ as \emph{effectively supported} on a ball of radius $R(t)\sqrt t$,
with an error that is smaller than the desired SFTE remainder scale.
This is exactly what is needed to feed into the insertion RG map at scale $L_{j_t}\simeq \sqrt t$.

We now state the flow-localization theorem used in this submission.

\begin{theorem}[Flow quasi-locality and derivative bounds]\label{thm:flow_local_v91}
Fix $t_\ast>0$ and a finite derivative order $r_{\max}$.
Under the bounded-regime control of Proposition~\ref{ass:bounded_regime}
(verified on the good-flow chart regime by Lemma~\ref{lem:chart_persist} and absorbed in expectations by Corollary~\ref{cor:superpoly_flow})
and the Gaussian kernel bounds~\ref{ass:gaussian_kernel}
(proved in Appendix~\ref{app:kernel_parametrix} and Appendix~\ref{app:discrete_parametrix}),
the following quasi-local truncation estimate holds for $*\in\{S,Q\}$:
\begin{enumerate}[label=(L\arabic*)]
\item (\textbf{Quasi-locality}) The functional derivatives of $\mathcal O_{*,t}(x)$ with respect to the initial field at points $y$ obey Gaussian off-diagonal bounds of the form
\eqref{eq:der_OS}--\eqref{eq:der_OQ} for all $1\le r\le r_{\max}$.
\item (\textbf{Effective support}) For any $\eta>0$ there exists a radius function $R(t)\asymp \sqrt{\log(1/t)}$ such that
modifying initial data outside $B_{R(t)}(x,t)$ changes $\mathcal O_{*,t}(x)$ by at most $O(t^{\eta})$ in sup norm.
\item (\textbf{Uniformity}) All constants can be chosen independent of lattice spacing $a$ (for $t\ge c a^2$) provided the chart/boundedness hypothesis is uniform in $a$.
\end{enumerate}
\end{theorem}

\paragraph{Status in this submission.}
The bounded-regime/chart hypothesis~\ref{ass:bounded_regime} is verified on the good-flow chart event by
Lemma~\ref{lem:chart_persist} (Section~\ref{sec:good_flow_chart}), and its complement is absorbed in expectations by
Corollary~\ref{cor:superpoly_flow}.
The Gaussian kernel bounds~\ref{ass:gaussian_kernel} are proved in Appendix~\ref{app:kernel_parametrix} (continuum parametrix)
and Appendix~\ref{app:discrete_parametrix} (uniform discrete analogue).
With these ingredients in place, all subsequent uses of ``flow localization'' in the insertion RG are bookkeeping.


\subsection{Auxiliary: good-flow chart event and bounded-regime control}
\label{sec:good_flow_chart}
This module proves the correct \emph{quasi-local} (Gaussian off-diagonal) functional-derivative bounds for
\[
\mathcal O_{S,t}(x)=\tfrac14\mathrm{tr}\,G_{\mu\nu}(t,x)G_{\mu\nu}(t,x),
\qquad
\mathcal O_{Q,t}(x)=\tfrac14\mathrm{tr}\,G_{\mu\nu}(t,x)\widetilde G_{\mu\nu}(t,x),
\]
The required analytic inputs are proved within this manuscript:
\begin{enumerate}[label=(\roman*)]
\item the \textbf{bounded-regime/chart} control for the flowed gauge potential $B(t)$ is proved in Lemma~\ref{lem:chart_persist}
(starting from the initial chart estimates in Lemma~\ref{lem:tree_chart_local}), and
\item the \textbf{Gaussian kernel bounds} for the linearized uniformly-parabolic operator are proved in Appendix~\ref{app:kernel_parametrix}
(continuum parametrix) and Appendix~\ref{app:discrete_parametrix} (lattice parametrix).
\end{enumerate}

This module attacks item (i): it replaces the ad hoc bounded-regime assumption by an explicit \textbf{good-flow chart mechanism}
that is compatible with the project's patched construction and subsequent de-patching estimates.

\begin{enumerate}[label=(A\arabic*)]
\item \textbf{Deterministic good-flow regime:} a sufficient condition on initial local curvature that implies the flow stays in a fixed chart up to time $t_\ast$.
\item \textbf{Patched/truncated state:} define the relevant expectations using a \emph{state patched by chart truncation} at the RG level,
and show the difference from the unpatched expectation is superpolynomially small in the AF window and can be absorbed into A1 remainders.
\end{enumerate}

We use the intrinsic distance-to-identity function
\[
\theta(U):=\mathrm{dist}_{G}(U,\mathbf 1),
\]
where $\mathrm{dist}_{G}$ is the geodesic distance for the standard bi-invariant metric induced by
the fixed bi-invariant Riemannian metric induced by the chosen inner product on $\mathfrak g$.
On the injectivity-radius chart, if $U=\exp(X)$ with $\|X\|\le \pi$ (Hilbert--Schmidt norm), then $\theta(U)=\|X\|$.


Let $\Lambda\subset\mathbb Z^4$ be a finite set of plaquettes (e.g.\ those whose geometric support intersects a fixed continuum compact $K$ enlarged by radius $R\sqrt t$).

\begin{definition}[Local small-curvature event]\label{def:good_plaq}
Fix $\varepsilon\in(0,\pi/2]$. Define
\[
G_{\Lambda}(\varepsilon):=\bigcap_{p\in \Lambda}\{\theta(U_p)\le \varepsilon\}.
\]
\end{definition}

\paragraph{Remark.}
$G_{\Lambda}(\varepsilon)$ is gauge invariant and only depends on plaquette holonomies.

We use the DeTurck gauge-fixed flow for analysis. Let $B(t)$ be the DeTurck-flowed gauge potential,
with initial data extracted from link variables via the principal logarithm on a fixed chart.

\begin{proposition}[Initial chart and small-curvature regime (proved in Lemma~\ref{lem:tree_chart_local} and Lemma~\ref{lem:local_tail})]\label{ass:init_chart}
There exist $\varepsilon\in(0,\pi/4]$ and $\delta\in(0,\pi/2)$ such that:
\begin{enumerate}[label=(IC\arabic*)]
\item all links in the relevant neighborhood admit exponential coordinates $U_e=\exp(A_e)$ with $\|A_e\|\le \delta$;
\item the plaquette angles in that neighborhood satisfy $G_{\Lambda}(\varepsilon)$ for a suitable plaquette set $\Lambda$.
\end{enumerate}
\end{proposition}

\begin{lemma}[Tree-gauge chart from small plaquettes (local version)]\label{lem:tree_chart_local}
There exists a constant $C_{\mathrm{geo}}$ depending only on the local block geometry (not on $a$) such that:
on $G_{\Lambda}(\varepsilon)$, after fixing tree gauge on each unit block intersecting the neighborhood, all unfixed links satisfy
\[
\theta(U_e)\le C_{\mathrm{geo}}\varepsilon.
\]
Hence choosing $\varepsilon\le \delta/C_{\mathrm{geo}}$ guarantees $\|A_e\|\le \delta$ throughout the neighborhood.
\end{lemma}

\begin{lemma}[Short-time persistence of the chart]\label{lem:chart_persist}
Assume \ref{ass:init_chart}. Fix $t_\ast>0$ and assume the DeTurck flow is run on $[0,t_\ast]$.
Then there exists $\delta_\ast=\delta_\ast(\delta,\varepsilon,t_\ast)$ such that:
\begin{enumerate}[label=(P\arabic*)]
\item $B(t)$ remains in the exponential chart on $[0,t_\ast]$ with $\|B(t)\|_\infty\le \delta_\ast$;
\item spatial derivatives satisfy the parabolic smoothing bounds
\[
\|\nabla^m B(t)\|_\infty\ \le\ C_m(\delta_\ast)\,t^{-m/2},\qquad 1\le m\le m_{\max},
\]
for any fixed $m_{\max}$ required by the insertion seminorm.
\end{enumerate}
\end{lemma}

\begin{proof}
In DeTurck gauge the flow is a semilinear uniformly parabolic system with smooth coefficients on the chart domain.
Local existence and uniqueness is standard; for lattice discretization with compact gauge group $G$ it is global in time as an ODE system.
Parabolic maximum principle plus Gr\"onwall bounds yield $\|B(t)\|_\infty\le \delta_\ast$ on $[0,t_\ast]$.
Standard parabolic Schauder/$L^p$ estimates yield $\|\nabla^m B(t)\|_\infty\lesssim t^{-m/2}$ with constants controlled by $\delta_\ast$.
\end{proof}

\paragraph{Consequence for kernel bounds.}
Lemma~\ref{lem:chart_persist} supplies the coefficient boundedness needed to invoke Aronson-type Gaussian kernel bounds for the linearized operator,
with constants depending only on $(\delta_\ast,m_{\max})$, hence uniform in $a$ (provided the chart holds uniformly in $a$ on the neighborhood).

The deterministic path above is useful but does not explain why the hypotheses should hold \emph{typically} under the Euclidean state.
We now provide a patched/truncated construction analogous to the patch-and-absorb step used later in the multiscale RG construction.

Let $\langle\cdot\rangle$ denote the underlying Euclidean expectation produced by the continuum construction (OS positivity is invoked only on the unflowed positive-time algebra).
For each fixed compact $K\subset\{x_0\ge\varepsilon_0\}$ and each fixed flow time $t\in(0,t_\ast]$, choose a neighborhood $\Lambda=\Lambda(K,t,R)$ of plaquettes
covering the effective support region of size $R\sqrt t$ around $K$.

\begin{definition}[Truncated expectation for flowed local insertions]\label{def:trunc_expect_flow}
Fix $\varepsilon\in(0,\pi/4]$.
For any positive-time observable $F$ supported in the neighborhood corresponding to $\Lambda$, define
\[
\langle F\rangle^{\mathrm{good}}_{K,t}
:=\frac{\langle F\mathbf 1_{G_{\Lambda}(\varepsilon)}\rangle}{\langle \mathbf 1_{G_{\Lambda}(\varepsilon)}\rangle}.
\]
\end{definition}

\begin{lemma}[Truncation stability (Cauchy--Schwarz form)]\label{lem:trunc_CS}
If $|F|\le 1$ then
\[
\big|\langle F\rangle-\langle F\rangle^{\mathrm{good}}_{K,t}\big|
\ \le\ 2\,\sqrt{\mathbb P(G_{\Lambda}(\varepsilon)^c)},
\]
where $\mathbb P(\cdot)$ is probability under the underlying state.
\end{lemma}

\begin{proof}
Write $\langle F\rangle=\langle F\mathbf 1_G\rangle+\langle F\mathbf 1_{G^c}\rangle$ and use $|\langle F\mathbf 1_{G^c}\rangle|\le \langle \mathbf 1_{G^c}\rangle$.
Normalize by $\langle \mathbf 1_G\rangle$ and use $\langle\mathbf 1_G\rangle\ge 1-\langle\mathbf 1_{G^c}\rangle$.
A sharper bound uses Cauchy--Schwarz: $|\langle F\mathbf 1_{G^c}\rangle|\le \|F\|_2\sqrt{\mathbb P(G^c)}\le \sqrt{\mathbb P(G^c)}$.
\end{proof}

In earlier finite-range fluctuation settings, one-step fluctuation measures were product/finite-range, making $\mathbb P(G^c)$ exponentially small by a direct one-plaquette bound.
For the full gauge state, correlations exist; nevertheless, in the AF window one expects a comparable bound for local small-curvature events.

We recall the local plaquette tail bound from Lemma~\ref{lem:local_tail} (proved in Section~\ref{sec:patching}).

\begin{corollary}[Superpolynomial smallness of the truncation error]\label{cor:superpoly_flow}
Assume Lemma~\ref{lem:local_tail}. Fix $K$ and $t$. Then for every $N$ there exists $u_0(N)>0$ such that in the AF window,
\[
u\le u_0(N)\quad\Longrightarrow\quad
\mathbb P(G_{\Lambda}(\varepsilon)^c)\le u^{2N}
\quad\text{and hence}\quad
|\langle F\rangle-\langle F\rangle^{\mathrm{good}}_{K,t}|\le 2u^{N}.
\]
\end{corollary}

\paragraph{Absorption into A1.}
Since A1 remainders only need $O(t^\eta)$ control with $\eta\ge 1$, the $u^N$ truncation errors can be made smaller than $t^\eta$ for any fixed $t$
by choosing $N$ large and working sufficiently deep in the AF window. Thus truncation errors are absorbed into the A1 remainder term.

Combine:
\begin{remark}[How this module is used in the main proof]\MainPath
The ad hoc bounded-regime assumption is eliminated: the good-flow chart event $G_{\Lambda}(\varepsilon)$,
together with Lemma~\ref{lem:tree_chart_local}, the tail bound Lemma~\ref{lem:local_tail},
and the persistence Lemma~\ref{lem:chart_persist}, provides a deterministic chart regime on $[0,t_\ast]$ with explicit derivative bounds.
The Gaussian kernel bounds required for the insertion seminorm are proved in Appendix~\ref{app:kernel_parametrix}
(continuum) and Appendix~\ref{app:discrete_parametrix} (lattice).
\end{remark}


\subsection{Auxiliary: Gaussian kernel bounds for the linearized DeTurck operator}
Lane~A's flow-localization bounds for the probes $\cO_{S,t}$ and $\cO_{Q,t}$ require the Gaussian off-diagonal bounds
\eqref{eq:K_gauss}--\eqref{eq:K_der_gauss} for the fundamental solution of the linearized DeTurck evolution.

Under the bounded-regime hypothesis~\ref{ass:bounded_regime}, the DeTurck-flowed gauge potential remains inside a fixed exponential chart on $(0,t_\ast]$
and satisfies deterministic pointwise bounds on a finite set of local derivatives.
As a consequence, the coefficients of the linearized equation are bounded (with at most time-integrable singularities inherited from $\|\nabla B(t)\|_\infty\lesssim t^{-1/2}$).

\paragraph{Status in this submission.}
The bounded-regime hypothesis~\ref{ass:bounded_regime} is verified on the good-flow chart event by Lemma~\ref{lem:chart_persist}
(Section~\ref{sec:good_flow_chart}); the contribution of the complement is absorbed in expectations by Corollary~\ref{cor:superpoly_flow}.

For the parametrix method, it is convenient to rewrite first-order terms in divergence form.

\begin{lemma}[Divergence-form reduction and coefficient bounds]\label{lem:div_form_reduction}
Let $\psi$ solve the linearized DeTurck system
\[
\partial_t \psi
=
\Delta \psi + \sum_{\alpha=1}^4 A_\alpha(t,x)\,\partial_\alpha\psi + V(t,x)\,\psi
\]
on $0<s<t\le t_\ast$, where $A_\alpha,V\in \mathrm{End}(V)$.
Assume that under the bounded-regime hypothesis~\ref{ass:bounded_regime} the coefficients obey
\[
\|A_\alpha(t,\cdot)\|_\infty\le \mathsf A_0,\qquad
\|V(t,\cdot)\|_\infty\le \mathsf V_0+\mathsf V_1\,t^{-1/2},
\qquad
\|\partial_\alpha A_\alpha(t,\cdot)\|_\infty\le \mathsf A_1\,t^{-1/2}.
\]
Then $\psi$ also satisfies the divergence-form system
\begin{equation}\label{eq:div_form_lin}
\partial_t \psi
=
\Delta \psi + \sum_{\alpha=1}^4 \partial_\alpha\big(A_\alpha(t,x)\psi\big)+W(t,x)\psi,
\end{equation}
with
\[
W(t,x):=V(t,x)-\sum_{\alpha=1}^4\partial_\alpha A_\alpha(t,x),
\qquad
\|W(t,\cdot)\|_\infty\le \mathsf W_0+\mathsf W_1\,t^{-1/2},
\]
where $\mathsf W_0=\mathsf V_0$ and $\mathsf W_1=\mathsf V_1+4\mathsf A_1$.
\end{lemma}

\begin{proof}
Use $A_\alpha\partial_\alpha\psi=\partial_\alpha(A_\alpha\psi)-(\partial_\alpha A_\alpha)\psi$ and regroup terms.
\end{proof}

\paragraph{Comment.}
In the concrete setting, $A_\alpha$ and $V$ are smooth functions of the background field $B(t)$ and its first derivatives
(DeTurck gauge).
Under Proposition~\ref{ass:bounded_regime}, $\|B(t)\|_\infty\le M_0$ and $\|\nabla B(t)\|_\infty\le M_1 t^{-1/2}$ on $(0,t_\ast]$,
so the bounds in Lemma~\ref{lem:div_form_reduction} hold with explicit $\mathsf A_0,\mathsf V_0,\mathsf V_1,\mathsf A_1$
depending on $(M_0,M_1)$ and universal Lie/gauge constants.

\begin{theorem}[Gaussian kernel bounds for \eqref{eq:div_form_lin}]\label{thm:gauss_continuum_parametrix}
Assume the hypotheses of Lemma~\ref{lem:div_form_reduction}.
Let $K(t,s;x,y)$ be the fundamental solution (matrix kernel) of \eqref{eq:div_form_lin}.
Then there exist constants $c_1\in(0,1)$ and, for each $m\in\{0,1,\dots,m_{\max}\}$, constants $C_m<\infty$ depending only on
\[
(d=4,\dim V,\mathsf A_0,\mathsf W_0,\mathsf W_1,t_\ast,m_{\max}),
\]
such that for all $0<s<t\le t_\ast$ and all $x,y\in\mathbb R^4$,
\begin{align}
\|K(t,s;x,y)\|_{\mathrm{op}}
&\le
C_0\,(t-s)^{-2}\exp\!\Big(-c_1\,\frac{|x-y|^2}{t-s}\Big),\\
\|\nabla_x^m K(t,s;x,y)\|_{\mathrm{op}}
&\le
C_m\,(t-s)^{-(2+m/2)}\exp\!\Big(-c_1\,\frac{|x-y|^2}{t-s}\Big).
\end{align}
\end{theorem}

\begin{proof}
This is proved by a Picard/parametrix expansion around the heat kernel, with explicit control of the Volterra majorant generated by
the $(t-r)^{-1/2}$ singularity coming from the divergence-form first-order term.
A fully written proof with constant dependence is given in Appendix~\ref{app:kernel_parametrix}
(Theorem~\ref{thm:gauss_cont_proved}); the bounds above are the specialization to $d=4$ with the exponential dependence
absorbed into the constants $C_m$ for $t\le t_\ast$.
\end{proof}

On $a\mathbb Z^4$ the DeTurck flow is implemented with finite-range difference operators.
The linearized evolution takes the form
\[
\partial_t\psi
=
\Delta_a\psi+\sum_{\alpha=1}^4\nabla_{a,\alpha}\big(A_{\alpha,a}(t,x)\psi\big)+W_a(t,x)\psi,
\]
with the same uniform coefficient bounds as in Lemma~\ref{lem:div_form_reduction} under Proposition~\ref{ass:bounded_regime}.

Once one inputs the uniform lattice heat-kernel bounds proved internally in Appendix~\ref{app:discrete_heat_kernel}
(Theorem~\ref{thm:lattice_gauss_uniform}), the lattice parametrix proof in Appendix~\ref{app:discrete_parametrix} yields
\eqref{eq:K_gauss}--\eqref{eq:K_der_gauss} with constants independent of $a$ (for $t-s\ge a^2$ and $x,y$ in the compact support region of the observable family).
In particular, the ``discrete input'' used here is now fully closed within this submission.

Under the bounded-regime hypothesis~\ref{ass:bounded_regime}, Lemma~\ref{lem:div_form_reduction} supplies the coefficient bounds required by the parametrix theorem.
Therefore Proposition~\ref{ass:gaussian_kernel} holds (in the continuum and, with the discrete heat kernel input, uniformly in $a$),
closing the flow-localization lemma stack needed for Theorem~\ref{thm:flow_local_v91}.


\subsection{Auxiliary: insertion RG and OS-seminorm bounds}
The Lane~A1 lemma stack isolates the \emph{insertion RG map} as a distinct analytic burden.
The main output of this subsection is the SFTE-style decomposition theorem (Theorem~\ref{thm:ins_decomp}), which decomposes a flowed insertion into finitely many renormalized local operators plus an extracted-irrelevant remainder that is small as $t\downarrow 0$.

It interfaces explicitly with the chart truncation / de-patching mechanism of Section~\ref{sec:patching} and with the chart-truncated integration philosophy of Definitions~\ref{def:trunc_measure} and~\ref{def:RG_good}.

\paragraph{What this subsection does.}

\begin{enumerate}[label=(\roman*)]
\item Define the Banach spaces of \emph{bulk} polymer activities $K_j$ and \emph{insertion} polymer activities $J_j$ at each dyadic scale.
\item Define a \textbf{patched} one-step insertion integration operator $\mathcal E^{\mathrm{good}}_{j+1}(u_j,K_j)$ that uses the good-block (chart) truncation and,
when flowed insertions are present, a local \emph{good-flow chart} truncation in the effective support region. On this event the deterministic hypotheses
for flow quasi-locality and kernel bounds hold by Lemma~\ref{lem:chart_persist} (Section~\ref{sec:good_flow_chart}).
\item Define the one-step \textbf{insertion RG map} $\mathcal R^{\mathrm{ins,good}}_j$ (integration + reblocking/rescaling + extraction).
\item Prove the \textbf{extracted-irrelevant contraction} and the resulting \textbf{SFTE-style decomposition} after iterating to a matching scale $\mu$,
with constants tracked.
\end{enumerate}

Fix dyadic scale factor $L=2$ and lattice spacing $a>0$. A \emph{$j$-block} is a cube of side $L^j a$.
A \emph{$j$-polymer} $X$ is a finite connected union of $j$-blocks. Let $|X|_j$ be the number of $j$-blocks in $X$.

\paragraph{Positive-time support.}
All insertions are smeared against test functions supported in a compact $K\subset\{x_0\ge \varepsilon\}$, so OS positivity is preserved.

We assume the chart-truncated dyadic RG map $\mathcal R^{\mathrm{good}}_j$ of Definition~\ref{def:RG_good} is available:
\[
(u_j,K_j)\ \mapsto\ (u_{j+1},K_{j+1}).
\]
We do \emph{not} reprove bulk RG bounds here; we only require the same seminorm infrastructure.

\begin{proposition}[Bulk RG seminorm axioms (proved in Proposition~\ref{ass:seminorm_axioms} and Proposition~\ref{ass:fluctuation_bounds})]\label{ass:bulk_axioms}
There exist scale-independent constants such that:
\begin{enumerate}[label=(B\arabic*)]
\item (\textbf{Product bound}) $\|FG\|_j\le C_{\mathrm{prod}}\|F\|_j\|G\|_j$ for local functionals on a $j$-polymer.
\item (\textbf{Counting}) polymer sums satisfy a bound with constant $C_{\mathrm{count}}$.
\item (\textbf{Integration bound}) the chart-truncated fluctuation expectation at step $j\to j+1$ defines a bounded linear operator on polymer activities with norm $\le C_E$.
\item (\textbf{Extraction}) projection onto relevant/marginal directions is bounded with constant $C_{\mathrm{ext}}$.
\end{enumerate}
\end{proposition}

Fix a symmetry sector and a finite operator basis $\{\mathcal O_i^{\mathrm{ren}}\}_{i=1}^M$ of gauge-invariant local operators
with engineering dimensions $\Delta_i$ and fixed quantum numbers (e.g.\ CP-even scalar, CP-odd pseudoscalar).

An \emph{insertion activity} at scale $j$ is a family
\[
J_j=\{J_j(X;f)\}_{X\in\mathscr P_j},
\]
where $\mathscr P_j$ is the set of $j$-polymers, each $J_j(X;f)$ is a local functional supported in $X$, and $J_j$ is \emph{linear} in the test function $f$.
We restrict to insertion activities supported on polymers that intersect the $j$-neighborhood of $\mathrm{supp}(f)$.

Fix $(k,K,\varepsilon)$ as in the OS seminorm. Define the insertion seminorm at scale $j$ by
\begin{equation}\label{eq:ins_norm}
\|J_j\|_{\mathrm{ins};j;k,K,\varepsilon}
:=
\sup_{\substack{f:\ \mathrm{supp}(f)\subset K\\ |f|_{C^k(K)}\le 1}}
\ \sup_{X\in\mathscr P_j}\ 
\Gamma_j(X)^{-1}\ \|J_j(X;f)\|_j,
\end{equation}
where $\|\cdot\|_j$ is the bulk local functional seminorm and $\Gamma_j(X)$ is the standard polymer weight (e.g.\ $\Gamma_j(X)=A^{|X|_j}$ with $A>1$).
All constants are required to be uniform in $a$ on the AF window.

\paragraph{Remark.}
This norm is the insertion analogue of the bulk polymer norm.
The additional supremum over $f$ integrates directly with the OS-seminorm topology used in Lemma~\ref{lem:ins_to_OS} and ultimately Theorem~\ref{thm:LaneA_closure}.

Fix an \emph{extraction depth} $d_\ast\in\{4,6\}$ (recommended: $d_\ast=6$ for safety).
Let $\mathcal V^{\le d_\ast}$ be the finite-dimensional span of basis operators with $\Delta_i\le d_\ast$ in the chosen sector.

At each scale $j$ choose normalized insertion basis activities $\Phi_{i,j}$ representing $\mathcal O_i^{\mathrm{ren}}$ at scale $j$
(e.g.\ block-localized versions with $\|\Phi_{i,j}\|_{\mathrm{ins};j}=1$).

\begin{definition}[Insertion extraction operator]\label{def:Ext_ins}
Define $\mathrm{Ext}^{\mathrm{ins}}_j$ as the bounded linear projection
\[
\mathrm{Ext}^{\mathrm{ins}}_j:\ J_j\ \mapsto\ \Big(\alpha_j,\ I_j\Big),
\qquad
J_j=\sum_{\Delta_i\le d_\ast}\alpha_{j,i}\,\Phi_{i,j}\ +\ I_j,
\]
where $I_j$ lies in the extracted-irrelevant insertion subspace (no components in $\mathcal V^{\le d_\ast}$).
\end{definition}

\begin{proposition}[Scaling gap after extraction (proved in Lemma~\ref{lem:Delta_min})]\label{ass:gap_ins}
Let $\Delta_{\min}=d_{\mathrm{next}}-d_\ast$, where $d_{\mathrm{next}}$ is the minimal engineering dimension of the next gauge-invariant operator in the sector.
Then under dyadic rescaling (block size doubled),
\[
\|\mathrm{Rescale}(I_j)\|_{\mathrm{ins};j+1}\ \le\ 2^{-\Delta_{\min}}\ \|I_j\|_{\mathrm{ins};j}.
\]
For pure gauge sectors: $d_\ast=4$ gives $\Delta_{\min}=2$; $d_\ast=6$ gives $\Delta_{\min}=4$.
\end{proposition}

At each step $j\to j+1$, the bulk RG uses a chart-truncated fluctuation expectation.
For insertions involving flowed fields, we additionally enforce a local good-flow chart event $G_{\Lambda}(\varepsilon)$
(Definition~\ref{def:good_plaq}) around the effective support region. On this event the bounded-regime bound \eqref{eq:bounds_B} holds
by Lemma~\ref{lem:chart_persist}, hence the deterministic hypotheses of Theorem~\ref{thm:flow_local_v91} and Proposition~\ref{ass:gaussian_kernel}
apply in the insertion neighborhood.

\begin{definition}[Patched expectation operator for insertion step]\label{def:E_good_ins}
Let $\mathbb E_{j+1}$ denote the original fluctuation expectation at step $j\to j+1$ and let $G_{j+1}$ denote the conjunction of:
(i) the good-block (small-curvature) chart event on each unit coarse block at this step, and
(ii) the local good-flow chart event $G_{\Lambda}(\varepsilon)$ in the insertion neighborhood (only needed when the insertion involves flowed fields).
Define
\[
\mathbb E^{\mathrm{good}}_{j+1}[F]
:=\frac{\mathbb E_{j+1}[F\,\mathbf 1_{G_{j+1}}]}{\mathbb E_{j+1}[\mathbf 1_{G_{j+1}}]}.
\]
\end{definition}

\paragraph{De-patching estimate.}
Lemma~\ref{lem:patch_error} bounds the difference between the unpatched and patched one-step insertion expectations in terms of $\mathbb P(G_{j+1}^c)$.
The explicit tail estimate for $\mathbb P(G_{j+1}^c)$ required in the main proof is provided by Lemma~\ref{lem:local_tail} (via the plaquette threshold choice defining $G_{j+1}$).

Let $Z_j(\varphi):=e^{-V_j(\varphi)}(1+K_j(\varphi))$ denote the bulk effective density at scale $j$ (in polymer form).
For an insertion activity $J_j$, define the \emph{pre-rescaled} insertion after integrating the $(j+1)$-fluctuation field $\zeta$ by the ratio
\begin{equation}\label{eq:ins_integration_ratio}
\widetilde J_{j+1}(\varphi)
:=
\frac{\mathbb E_{j+1}^{\mathrm{good}}\!\big[J_j(\varphi+\zeta)\,Z_j(\varphi+\zeta)\big]}
{\mathbb E_{j+1}^{\mathrm{good}}\!\big[Z_j(\varphi+\zeta)\big]}.
\end{equation}
In polymer coordinates, $\widetilde J_{j+1}$ is re-expanded as a $(j+1)$-scale insertion activity
$\widetilde J_{j+1}=\{\widetilde J_{j+1}(Y;f)\}_{Y\in\mathscr P_{j+1}}$.

\begin{definition}[One-step insertion integration operator]\label{def:E_ins}
Define the (bulk-dependent) linear map
\[
\mathcal E^{\mathrm{good}}_{j+1}(u_j,K_j):\ J_j\ \mapsto\ \widetilde J_{j+1}
\]
by \eqref{eq:ins_integration_ratio} together with the standard polymer re-expansion (same combinatorics as in the bulk RG).
\end{definition}

\begin{lemma}[Boundedness of $\mathcal E^{\mathrm{good}}_{j+1}$ in insertion seminorm]\label{lem:Eins_bound}
Assume the bulk RG ball conditions (small $u_j$ and small $\|K_j\|_j$) and the chart-truncated fluctuation bounds of Proposition~\ref{ass:fluctuation_bounds}.
Then there exists a scale-independent constant $C_{E,\mathrm{ins}}$ such that
\[
\|\mathcal E^{\mathrm{good}}_{j+1}(u_j,K_j)J_j\|_{\mathrm{ins};j+1}
\le C_{E,\mathrm{ins}}\ \|J_j\|_{\mathrm{ins};j},
\]
uniformly in $a$ and uniformly in $j$.
\end{lemma}

\paragraph{Comment (flowed insertions).}
When the insertion involves flowed fields, the patched expectation in Definition~\ref{def:E_good_ins} includes the local good-flow chart event.
On this event the deterministic hypotheses needed for the flow quasi-locality and kernel bounds are satisfied by Lemma~\ref{lem:chart_persist}
(Section~\ref{sec:good_flow_chart}); the contribution of the complement is absorbed (at fixed $t$) by the truncation estimate
Corollary~\ref{cor:superpoly_flow} together with Lemma~\ref{lem:patch_error}.

\paragraph{Proof.}
The proof is given in Proposition~\ref{prop:Eins_bound} (see the next subsection).

\begin{definition}[Patched one-step insertion RG map]\label{def:RG_ins}
Given $(u_j,K_j)$ and an insertion activity $J_j$, define:
\begin{enumerate}[label=(\arabic*)]
\item integrate: $\widetilde J_{j+1}=\mathcal E^{\mathrm{good}}_{j+1}(u_j,K_j)J_j$;
\item extract: $(\alpha_{j+1},\widehat I_{j+1})=\mathrm{Ext}^{\mathrm{ins}}_{j+1}(\widetilde J_{j+1})$;
\item rescale the remainder: $I_{j+1}=\mathrm{Rescale}(\widehat I_{j+1})$.
\end{enumerate}
Define the insertion RG map
\[
\mathcal R^{\mathrm{ins,good}}_j:\ (u_j,K_j;\alpha_j,I_j)\ \longmapsto\ (u_{j+1},K_{j+1};\alpha_{j+1},I_{j+1}),
\]
where $(u_{j+1},K_{j+1})$ is the bulk RG image under $\mathcal R^{\mathrm{good}}_j$.
\end{definition}

\begin{proposition}[One-step contraction on extracted-irrelevant insertion sector]\label{prop:ins_contract}
Under Propositions~\ref{ass:bulk_axioms} and~\ref{ass:gap_ins}, and Lemma~\ref{lem:Eins_bound},
the remainder map satisfies
\[
\|I_{j+1}\|_{\mathrm{ins};j+1}
\le \kappa_{\mathrm{ins}}\ \|I_j\|_{\mathrm{ins};j},
\qquad
\kappa_{\mathrm{ins}}:=C_{E,\mathrm{ins}}\,2^{-\Delta_{\min}}.
\]
In particular, if $\kappa_{\mathrm{ins}}<1$ (achieved either by bounding $C_{E,\mathrm{ins}}$ or by choosing deeper extraction to increase $\Delta_{\min}$),
then the extracted-irrelevant sector contracts uniformly in $j$ and uniformly in $a$.
\end{proposition}

Fix a renormalization scale $\mu>0$ and a flow time $t\in(0,t_\ast]$ with $t\gg a^2$.
Let $n(t,\mu)$ be the number of dyadic steps needed so that the block size $L^{n}\sqrt t$ is comparable to $\mu^{-1}$:
\begin{equation}\label{eq:n_def}
n(t,\mu):=\Big\lceil \log_2\Big(\frac{1}{\mu\sqrt t}\Big)\Big\rceil.
\end{equation}

\begin{theorem}[Insertion RG decomposition at scale $\mu$]\label{thm:ins_decomp}
Assume $\kappa_{\mathrm{ins}}<1$ and iterate the patched insertion RG map $\mathcal R^{\mathrm{ins,good}}$ for $n=n(t,\mu)$ steps.
Then the flowed insertion $\mathcal O_{*,t}(f)$ admits a decomposition
\[
\mathcal O_{*,t}(f)=\sum_{\Delta_i\le d_\ast}\alpha_i(t,\mu)\,\mathcal O_i^{\mathrm{ren}}(f;\mu)\ +\ \mathcal R_{*,t}(f),
\]
where $\alpha_i(t,\mu)$ are the extracted coefficients at the matching scale and the remainder satisfies the insertion norm bound
\begin{equation}\label{eq:rem_bound}
\|\mathcal R_{*,t}\|_{\mathrm{ins};\mu;k,K,\varepsilon}
\le C_{\mathrm{init}}\ \kappa_{\mathrm{ins}}^{\,n(t,\mu)}.
\end{equation}
Consequently,
\begin{equation}\label{eq:rem_tpower}
\|\mathcal R_{*,t}\|_{\mathrm{ins};\mu;k,K,\varepsilon}
\le C_{\mathrm{init}}\,C_\mu\ (\mu^2 t)^{\eta},
\qquad
\eta:=\frac{\Delta_{\min}-\log_2 C_{E,\mathrm{ins}}}{2}\ >\ 0,
\end{equation}
with $C_\mu$ depending only on $\mu$ and fixed dyadic rounding constants.
\end{theorem}

\begin{proof}
Iterate Proposition~\ref{prop:ins_contract} to obtain \eqref{eq:rem_bound}.
Write $\kappa_{\mathrm{ins}}=C_{E,\mathrm{ins}}2^{-\Delta_{\min}}$ and use $2^{n(t,\mu)}\asymp 1/(\mu\sqrt t)$ to bound
$C_{E,\mathrm{ins}}^{n}\le (2^{n})^{\log_2 C_{E,\mathrm{ins}}}\lesssim (\mu\sqrt t)^{-\log_2 C_{E,\mathrm{ins}}}$ and
$2^{-\Delta_{\min}n}\lesssim (\mu\sqrt t)^{\Delta_{\min}}$.
Combine to obtain \eqref{eq:rem_tpower}.
\end{proof}

\paragraph{Interpretation.}
The effective exponent $\eta$ is the engineering scaling gap $\Delta_{\min}/2$ \emph{reduced} by the per-step integration-norm growth $\log_2 C_{E,\mathrm{ins}}/2$.
Any $\eta>0$ suffices for the Lane~A closure theorem (Theorem~\ref{thm:LaneA_closure}). If one wants $\eta\ge 1$, choose extraction depth $d_\ast$ large enough (e.g.\ $d_\ast=6$ gives $\Delta_{\min}=4$)
and/or improve the bound on $C_{E,\mathrm{ins}}$.

Define the unpatched insertion integration operator $\mathcal E_{j+1}$ by replacing $\mathbb E^{\mathrm{good}}_{j+1}$ with $\mathbb E_{j+1}$ in \eqref{eq:ins_integration_ratio}.
The difference is controlled by the probability of the complement of $G_{j+1}$.

\begin{lemma}[Patched vs.\ unpatched one-step insertion expectations]\label{lem:patch_error}
Let $G_{j+1}$ be any measurable ``good'' event and define the patched expectation
\[
\mathbb E^{\mathrm{good}}_{j+1}[F]
:=\frac{\mathbb E_{j+1}[F\,\mathbf 1_{G_{j+1}}]}{\mathbb E_{j+1}[\mathbf 1_{G_{j+1}}]}.
\]
Then for any bounded local insertion functional $F$,
\begin{equation}\label{eq:patch_error_general}
\big|\mathbb E_{j+1}[F]-\mathbb E^{\mathrm{good}}_{j+1}[F]\big|
\ \le\ 2\|F\|_\infty\,\sqrt{\mathbb P(G_{j+1}^c)}.
\end{equation}
In particular, choosing $G_{j+1}$ as a local plaquette-smallness event and applying Lemma~\ref{lem:local_tail}
yields an explicit superpolynomially small error in the AF window (absorbed into the Lane~A remainder).
\end{lemma}

\begin{proof}
Write
\[
\mathbb E_{j+1}[F]-\mathbb E^{\mathrm{good}}_{j+1}[F]
=
\frac{\mathbb E_{j+1}[F\,\mathbf 1_{G_{j+1}^c}]}{\mathbb E_{j+1}[\mathbf 1_{G_{j+1}}]}
-
\frac{\mathbb E_{j+1}[F]\,\mathbb P(G_{j+1}^c)}{\mathbb E_{j+1}[\mathbf 1_{G_{j+1}}]}.
\]
Using $\mathbb E_{j+1}[\mathbf 1_{G_{j+1}}]\ge 1-\mathbb P(G_{j+1}^c)\ge 1/2$ when $\mathbb P(G_{j+1}^c)\le 1/2$ (which is the regime of interest),
and Cauchy--Schwarz gives $|\mathbb E_{j+1}[F\,\mathbf 1_{G_{j+1}^c}]|\le \|F\|_\infty\sqrt{\mathbb P(G_{j+1}^c)}$.
The bound \eqref{eq:patch_error_general} follows.
\end{proof}

\begin{itemize}
\item The insertion RG map $\mathcal R^{\mathrm{ins,good}}$ is defined in the same dyadic seminorm framework as the bulk RG, and its interface to the local bounded-flow truncation (used only for flowed insertions) is explicit.
\item The extracted-irrelevant contraction depends on two explicit constants: the scaling gap $\Delta_{\min}$ (from the bulk RG) and the insertion integration bound $C_{E,\mathrm{ins}}$, bounded uniformly by Proposition~\ref{prop:Eins_bound}.
\item The de-patching error between patched and unpatched expectations is controlled by Lemma~\ref{lem:patch_error} together with the explicit tail estimate Lemma~\ref{lem:local_tail}.
\end{itemize}

In the previous subsection we defined the good-event (chart-truncated) insertion RG map $\mathcal R^{\mathrm{ins,good}}$ and reduced the A1 ``insertion RG'' burden to two explicit constants:
the scaling gap $\Delta_{\min}$ (already fixed by extraction depth) and the one-step insertion integration constant $C_{E,\mathrm{ins}}$.

We \textbf{prove} the one-step boundedness lemma stated below.

\medskip
\noindent\textbf{Lemma~\ref{lem:Eins_bound} (proved here).}
There exists a scale-independent constant $C_{E,\mathrm{ins}}$ such that on the RG ball,
\[
\|\mathcal E^{\mathrm{good}}_{j+1}(u_j,K_j)J_j\|_{\mathrm{ins};j+1}\le C_{E,\mathrm{ins}}\|J_j\|_{\mathrm{ins};j},
\]
uniformly in dyadic scale $j$ and uniformly in $a\to0$.

\medskip
We give an explicit formula for $C_{E,\mathrm{ins}}$ in terms of the \emph{already-defined} constants and the RG ball radius of $K$.

At scale $j$ the bulk effective density is
\[
Z_j(\varphi)=e^{-V_j(\varphi)}\big(1+K_j(\varphi)\big),
\]
and the good-event fluctuation expectation at step $j\to j+1$ is $\mathbb E^{\mathrm{good}}_{j+1}$ (Definition~\ref{def:E_good_ins}).

Given an insertion activity $J_j$, define the pre-rescaled integrated insertion by the ratio
\begin{equation}\label{eq:ratio}
\widetilde J_{j+1}(\varphi)
:=
\frac{\mathbb E^{\mathrm{good}}_{j+1}\!\big[J_j(\varphi+\zeta)\,Z_j(\varphi+\zeta)\big]}
{\mathbb E^{\mathrm{good}}_{j+1}\!\big[Z_j(\varphi+\zeta)\big]}.
\end{equation}
The one-step insertion integration operator is $\mathcal E^{\mathrm{good}}_{j+1}(u_j,K_j):J_j\mapsto \widetilde J_{j+1}$ followed by the standard polymer re-expansion.

We adopt the same local functional seminorm $\|\cdot\|_j$ and polymer weights $\Gamma_j(X)$ used in Proposition~\ref{ass:seminorm_axioms}.
The insertion norm is ( \eqref{eq:ins_norm}):
\[
\|J_j\|_{\mathrm{ins};j;k,K,\varepsilon}
=\sup_{f}\sup_{X}\Gamma_j(X)^{-1}\|J_j(X;f)\|_j,
\]
where $f$ ranges over $\mathrm{supp}(f)\subset K$, $|f|_{C^k}\le 1$.

\begin{proposition}[Integration and counting constants (proved in Proposition~\ref{ass:seminorm_axioms} and Section~\ref{sec:patching})]\label{ass:C_constants}
There exist scale-independent constants:
\begin{itemize}
\item $C_E$ such that the \emph{bulk} one-step integration operator under the patched expectation is bounded on polymer activities:
\[
\|\mathcal E^{\mathrm{good}}_{j+1}(u_j,K_j)F_j\|_{j+1}\le C_E\|F_j\|_j,
\]
for all local functionals/polymers $F_j$ in the RG ball.
\item $C_{\mathrm{prod}}$ such that $\|FG\|_j\le C_{\mathrm{prod}}\|F\|_j\|G\|_j$ (product bound).
\item $C_{\mathrm{mark}}$ controlling the polymer re-expansion with \emph{one marked polymer} (insertion):
the number of ways to re-expand a marked polymer configuration into a $(j+1)$-polymer is bounded by $C_{\mathrm{mark}}$
relative to the unmarked case. One may take $C_{\mathrm{mark}}\le C_{\mathrm{count}}$ where $C_{\mathrm{count}}$ is the bulk counting constant.
\end{itemize}
All constants are uniform in $a\to0$ because they are established under the good-chart (chart-truncated) regime.
\end{proposition}

\paragraph{RG ball radius.}
Let $r_K$ be the uniform bound on the bulk irrelevant activity:
\[
\|K_j\|_j\le r_K\qquad\text{for all }j \text{ in the AF window}.
\]

The ratio \eqref{eq:ratio} can be reduced to expectations under the \emph{convex} reference weight $e^{-V_j}$, treating $K_j$ as a perturbation.
Define the normalized patched expectation with respect to $e^{-V_j}$:
\begin{equation}\label{eq:EVdef}
\mathbb E^{V,\mathrm{good}}_{j+1}[F]
:=
\frac{\mathbb E^{\mathrm{good}}_{j+1}\!\big[F(\varphi+\zeta)\,e^{-V_j(\varphi+\zeta)}\big]}
{\mathbb E^{\mathrm{good}}_{j+1}\!\big[e^{-V_j(\varphi+\zeta)}\big]}.
\end{equation}
(This is expectation under the patched fluctuation measure reweighted by $e^{-V_j}$; on the chart regime $V_j$ is uniformly convex by Theorem~\ref{thm:uniform_convexity}.)

\begin{lemma}[Ratio decomposition]\label{lem:ratio_decomp}
For each fixed coarse field $\varphi$,
\begin{equation}\label{eq:ratio_decomp}
\widetilde J_{j+1}(\varphi)
=
\frac{\mathbb E^{V,\mathrm{good}}_{j+1}[J_j]+\mathbb E^{V,\mathrm{good}}_{j+1}[J_j K_j]}
{1+\mathbb E^{V,\mathrm{good}}_{j+1}[K_j]}.
\end{equation}
Equivalently,
\[
\widetilde J_{j+1}=\mathbb E^{V,\mathrm{good}}_{j+1}[J_j]
+
\frac{\mathrm{Cov}^{V,\mathrm{good}}_{j+1}(J_j,K_j)}{1+\mathbb E^{V,\mathrm{good}}_{j+1}[K_j]},
\]
where $\mathrm{Cov}(J,K)=\mathbb E[JK]-\mathbb E[J]\mathbb E[K]$.
\end{lemma}

\begin{proof}
Factor $\mathbb E^{\mathrm{good}}_{j+1}[e^{-V_j}]$ out of numerator and denominator of \eqref{eq:ratio} and divide by it, yielding \eqref{eq:ratio_decomp}.
The covariance form is obtained by adding and subtracting $\mathbb E[J]\mathbb E[K]$.
\end{proof}

\begin{lemma}[Bound on $\mathbb E^{V,\mathrm{good}}_{j+1}[K_j$] and denominator lower bound]\label{lem:denom}
Under Proposition~\ref{ass:C_constants}, for all $j$ in the RG ball,
\[
\big\|\mathbb E^{V,\mathrm{good}}_{j+1}[K_j]\big\|_{j+1}\le C_E\,\|K_j\|_j\le C_E r_K.
\]
In particular, if $C_E r_K\le \tfrac12$, then
\begin{equation}\label{eq:denom_lower}
\Big\|\big(1+\mathbb E^{V,\mathrm{good}}_{j+1}[K_j]\big)^{-1}\Big\|_{j+1}\le \frac{1}{1-C_E r_K}\ \le\ 2.
\end{equation}
\end{lemma}

\begin{proof}
$\mathbb E^{V,\mathrm{good}}_{j+1}$ is a normalized version of the patched expectation; boundedness in the local seminorm is inherited from the bulk integration bound $C_E$.
The inverse bound follows from the Neumann series for $(1+X)^{-1}$ when $\|X\|<1$.
\end{proof}

\begin{lemma}[Local bound for the integrated insertion]\label{lem:local_bound}
Assume \ref{ass:C_constants} and $C_E r_K\le \tfrac12$. Then for any local insertion functional $J$ supported in a $j$-polymer $X$,
\[
\|\widetilde J_{j+1}\|_{j+1}
\le
\frac{C_E}{1-C_E r_K}\Big(\|J\|_j + C_{\mathrm{prod}}\|J\|_j\|K_j\|_j\Big)
\ \le\ 2C_E(1+C_{\mathrm{prod}}r_K)\|J\|_j.
\]
\end{lemma}

\begin{proof}
From \eqref{eq:ratio_decomp} and the product bound,
\[
\|\mathbb E^{V,\mathrm{good}}_{j+1}[J]\|_{j+1}\le C_E\|J\|_j,\qquad
\|\mathbb E^{V,\mathrm{good}}_{j+1}[J K_j]\|_{j+1}\le C_E\|J K_j\|_j\le C_E C_{\mathrm{prod}}\|J\|_j\|K_j\|_j.
\]
Multiply by the denominator inverse bound \eqref{eq:denom_lower}.
\end{proof}

We now lift the single-polymer bound to the insertion norm \eqref{eq:ins_norm}.

\begin{proposition}[Insertion integration bound (one-step)]\label{prop:Eins_bound}
Assume \ref{ass:C_constants} and $C_E r_K\le \tfrac12$.
Then the one-step insertion integration operator satisfies
\[
\|\mathcal E^{\mathrm{good}}_{j+1}(u_j,K_j)J_j\|_{\mathrm{ins};j+1;k,K,\varepsilon}
\le C_{E,\mathrm{ins}}\ \|J_j\|_{\mathrm{ins};j;k,K,\varepsilon},
\]
with an explicit scale-independent constant
\begin{equation}\label{eq:CEins}
C_{E,\mathrm{ins}}
:= C_{\mathrm{mark}}\ \frac{C_E\,(1+C_{\mathrm{prod}}r_K)}{1-C_E r_K}
\ \le\ 2C_{\mathrm{mark}}\,C_E\,(1+C_{\mathrm{prod}}r_K).
\end{equation}
All constants are uniform in $a\to0$ under the patched chart regime.
\end{proposition}

\begin{proof}
Expand $\widetilde J_{j+1}$ as a $(j+1)$-scale insertion activity by reblocking the marked polymer $X$ into the unique $(j+1)$-polymer $Y$ that contains it,
and summing over the finite-range neighborhoods required by the fluctuation integration. The number of such re-expansions is bounded by $C_{\mathrm{mark}}$.
Apply Lemma~\ref{lem:local_bound} on each marked configuration and absorb the combinatorial factor into $C_{\mathrm{mark}}$.
Finally take the supremum over $f$ in the $C^k$ unit ball; the constants do not depend on $f$.
\end{proof}

\begin{corollary}[Lemma~\ref{lem:Eins_bound} proved]\label{cor:lemEins}
Lemma~\ref{lem:Eins_bound} holds with $C_{E,\mathrm{ins}}$ given by \eqref{eq:CEins}.
\end{corollary}

With extraction depth $d_\ast=6$, pure-gauge sectors have $\Delta_{\min}=4$, hence the one-step irrelevant contraction constant is
\[
\kappa_{\mathrm{ins}}=C_{E,\mathrm{ins}}\,2^{-\Delta_{\min}}=\frac{C_{E,\mathrm{ins}}}{16}.
\]
Thus contraction holds if $C_{E,\mathrm{ins}}<16$.
Using \eqref{eq:CEins}, a sufficient condition is
\[
2C_{\mathrm{mark}}\,C_E\,(1+C_{\mathrm{prod}}r_K)\ <\ 16.
\]
This is an explicit finite-dimensional inequality once $(C_{\mathrm{mark}},C_E,C_{\mathrm{prod}},r_K)$ are instantiated from the bulk RG ball.
(No numerical experimentation is required: after the constants are fixed, the condition is a deterministic verification.)

\begin{theorem}[Canonical block one-shell increment]\label{thm:canonical_block_one_shell_increment}
Fix an exact-dimension symmetry block
\[
\mathfrak s=(\Delta,\epsilon,\lambda)
\]
of Definition~\ref{def:local_operator_spaces}, together with the canonical local basis, canonical flowed basis, and canonical block norm of
Appendix~\ref{app:normalization_crosscheck}.
Let
\[
\mathbf Y_{\mathfrak s,j}(f):=\mathbf Y_{\mathfrak s,t_j}(f)
\]
denote the vector of canonical flowed lifts at dyadic times $t_j=2^{-2j}t_0$, ordered in the fixed canonical basis.
Then there exist scale-independent constants $C_{\mathfrak s},C_{\mathfrak s,f}<\infty$ and, for each $j$, canonically defined linear maps
\[
A_{\mathfrak s,j}:\mathscr V_{\mathfrak s}\to \mathscr V_{\mathfrak s},
\qquad
B_{\mathfrak s\mathfrak s',j}:\mathscr V_{\mathfrak s'}\to \mathscr V_{\mathfrak s}
\quad (\mathfrak s'\prec \mathfrak s),
\]
such that
\begin{equation}\label{eq:canonical_block_one_shell}
\mathbf Y_{\mathfrak s,j+1}(f)
=
\bigl(I_{\mathfrak s}+u_jA_{\mathfrak s,j}\bigr)\mathbf Y_{\mathfrak s,j}(f)
+\sum_{\mathfrak s'\prec \mathfrak s}u_j B_{\mathfrak s\mathfrak s',j}\,\mathbf Y_{\mathfrak s',j}(f)
+\mathbf Q_{\mathfrak s,j}(f),
\end{equation}
with
\begin{equation}\label{eq:canonical_block_one_shell_bounds}
\|A_{\mathfrak s,j}\|_{\mathrm{op,can}}+\sum_{\mathfrak s'\prec \mathfrak s}\|B_{\mathfrak s\mathfrak s',j}\|_{\mathrm{op,can}}
\le C_{\mathfrak s},
\qquad
\|\mathbf Q_{\mathfrak s,j}(f)\|_{L^2(\|\cdot\|_{\mathrm{can}})}\le C_{\mathfrak s,f}\,u_j^2.
\end{equation}
Equivalently, after quotienting by the lower blocks of the grading order, the exact-dimension one-shell transport on $\mathscr V_{\mathfrak s}$ has the basis-free form
\[
S_j^{\mathfrak s}=I_{\mathfrak s}+u_jA_{\mathfrak s,j}+R_{\mathfrak s,j},
\qquad
\|R_{\mathfrak s,j}\|_{\mathrm{op,can}}\le C_{\mathfrak s}\,u_j^2.
\]
\end{theorem}

\begin{proof}
Fix one canonical generator in the block $\mathfrak s$.
The flow quasi-locality theorem (Theorem~\ref{thm:flow_local_v91}) and the patched insertion decomposition at matching scale (Theorem~\ref{thm:ins_decomp}) write the shell difference between $t_j$ and $t_{j+1}$ as a finite sum of local terms supported in an $O(\sqrt{t_j})$ neighborhood of the test-function support.
Because the local and flowed bases are written using the same complete-contraction formulas, every leading exact-dimension contribution lands back in the same block $\mathscr V_{\mathfrak s}$,
while terms carrying extra derivatives or explicit powers of $t_j$ land only in strictly lower blocks of the grading order.
This produces the linear maps $A_{\mathfrak s,j}$ and $B_{\mathfrak s\mathfrak s',j}$ in \eqref{eq:canonical_block_one_shell}.

The coefficient size is controlled uniformly on the AF ball by Proposition~\ref{prop:Eins_bound};
Theorem~\ref{thm:ins_to_OS} upgrades the insertion seminorm control to the OS/$L^2$ seminorm used later in Lane~A.
The remaining terms either contain at least two shell factors or one extracted-irrelevant insertion remainder,
and are therefore $O(u_j^2)$ in $L^2$ by Theorem~\ref{thm:ins_decomp} together with Proposition~\ref{prop:Eins_bound}.
Since the block basis and the canonical norm are fixed once and for all (Definition~\ref{def:LaneA_canonical_block_norm}), these bounds are operator bounds on the whole block rather than coordinatewise statements.
This yields \eqref{eq:canonical_block_one_shell_bounds}.
The quotient formulation is simply the projection of \eqref{eq:canonical_block_one_shell} to the exact-dimension block.
\end{proof}

\begin{corollary}[Graded triangular transport on every finite engineering cap]\label{cor:graded_triangular_transport_all_blocks}
Fix $\Delta_{\max}\ge 0$ and order all exact-dimension symmetry blocks with engineering dimension $\le \Delta_{\max}$ by the grading order used in Section~\ref{sec:sharp_local}.
Let
\[
\mathbf Y^{(\le \Delta_{\max})}_j(f)
\]
be the direct sum of their canonical flowed lifts.
Then there exist block-lower-triangular operators $A^{(\le \Delta_{\max})}_j$ and remainders $\mathbf Q^{(\le \Delta_{\max})}_j(f)$ such that
\begin{equation}\label{eq:finite_cap_one_shell}
\mathbf Y^{(\le \Delta_{\max})}_{j+1}(f)
=
\bigl(I+u_jA^{(\le \Delta_{\max})}_j\bigr)\mathbf Y^{(\le \Delta_{\max})}_j(f)
+\mathbf Q^{(\le \Delta_{\max})}_j(f),
\end{equation}
with
\[
\|A^{(\le \Delta_{\max})}_j\|_{\mathrm{op}}\le C_{\Delta_{\max}},
\qquad
\|\mathbf Q^{(\le \Delta_{\max})}_j(f)\|_{L^2}\le C_{\Delta_{\max},f}\,u_j^2.
\]
Moreover this block-lower-triangular transport is compatible under cap enlargement: if $\Delta'_{\max}\ge \Delta_{\max}$, the transport on the smaller cap is the restriction of that on the larger cap.
\end{corollary}

\begin{proof}
Apply Theorem~\ref{thm:canonical_block_one_shell_increment} block-by-block and assemble the finitely many blocks into a direct sum.
Lower-block mixing remains lower triangular after assembly, and the finite direct sum preserves the $O(u_j^2)$ remainder bound.
Compatibility under cap enlargement is exactly Lemma~\ref{lem:LaneA_block_cap_compat}.
\end{proof}

\begin{remark}[How the insertion bound is used after the basis-free hardening]
Proposition~\ref{prop:Eins_bound} now feeds two load-bearing steps.
First, together with Theorem~\ref{thm:ins_decomp} and Theorem~\ref{thm:ins_to_OS}, it proves the canonical block one-shell increment theorem above,
which is the unconditional input replacing the earlier phrase ``one-shell increment mechanism in each fixed finite operator basis'' in Theorem~\ref{thm:field_algebra_truncation}.
Second, it continues to control the off-diagonal blocks and remainders in the Lane~A inverse/closure theorems.
De-patching to the physical Wilson state is handled separately in Section~\ref{sec:patching} and does not alter $C_{E,\mathrm{ins}}$.
\end{remark}

We now compare two norm/topology controls used in the construction:
the dyadic insertion seminorm (a polymer/RG seminorm used to bound inserted observables scale-by-scale) and
the Osterwalder--Schrader (OS) seminorm used for OS reconstruction.
The key estimate is a one-way bound
\[
\|R_{*,t}\|_{\mathrm{OS};k,K,\varepsilon}\le C_{\mathrm{OS}\leftarrow\mathrm{ins}}\,
\|R_{*,t}\|_{\mathrm{ins};\mu;k,K,\varepsilon},
\]
showing that insertion-seminorm control implies OS-seminorm control for the same remainder observable $R_{*,t}$.
This is Lemma~\ref{lem:ins_to_OS} below.

\paragraph{What proves.}
Under the \emph{patched-chart} (good-block + good-flow) expectation state $\langle\cdot\rangle^{\mathrm{good}}$,
we prove a uniform constant $C_{\mathrm{OS}\leftarrow\mathrm{ins}}$ (independent of scale $j$ and lattice spacing $a$) such that for any
positive-time insertion activity $J$,
\[
\|F_J\|_{\mathrm{OS};k,K,\varepsilon}\ \le\ C_{\mathrm{OS}\leftarrow\mathrm{ins}}\ \|J\|_{\mathrm{ins};j;k,K,\varepsilon},
\]
where $F_J(f):=\sum_{X}J(X;f)$ is the associated observable.

\paragraph{Interpretation.}
The OS seminorm is controlled by a \emph{supremum} bound on $F_J$; under the patched-chart state the field space is compact on each relevant block,
so the local functional seminorm controls the sup norm.
Polymer weights and counting convert this local bound into a bound for the full summed observable.

Let $\langle\cdot\rangle^{\mathrm{good}}$ be the patched-chart (good-event, chart-truncated) expectation state $\langle\cdot\rangle^{\mathrm{good}}$.
Let $\Theta$ denote Euclidean time reflection and $\mathcal A_+$ the positive-time algebra.

For $F\in\mathcal A_+$ define
\[
\|F\|_{\mathrm{OS}}:=\sqrt{\langle \Theta F\cdot F\rangle^{\mathrm{good}}}.
\]
Fix $\varepsilon>0$ and compact $K\subset\{x_0\ge \varepsilon\}$ and define the OS seminorm (Lane~A topology)
\[
\|F\|_{\mathrm{OS};k,K,\varepsilon}
:=\sup\{\|F(f)\|_{\mathrm{OS}}:\ \mathrm{supp}(f)\subset K,\ |f|_{C^k(K)}\le 1\}.
\]

At dyadic scale $j$, an insertion activity is a family $J=\{J(X;f)\}_{X\in\mathscr P_j}$ supported on $j$-polymers.
The insertion seminorm is
\begin{equation}
\|J\|_{\mathrm{ins};j;k,K,\varepsilon}
:=
\sup_{\substack{f:\ \mathrm{supp}(f)\subset K\\ |f|_{C^k(K)}\le 1}}
\ \sup_{X\in\mathscr P_j}\ 
\Gamma_j(X)^{-1}\ \|J(X;f)\|_j,
\end{equation}
with polymer weight $\Gamma_j(X)=A^{|X|_j}$ for some fixed $A>1$ and local functional seminorm $\|\cdot\|_j$.

\begin{lemma}[OS bound by $L^\infty$]\label{lem:OS_by_sup}
For any $F\in\mathcal A_+$,
\[
\|F\|_{\mathrm{OS}}\le \|F\|_{L^\infty(\Omega)},
\]
where the sup is over the underlying configuration space $\Omega$ on which the patched probability measure is supported.
\end{lemma}

\begin{proof}
Since $\langle\cdot\rangle^{\mathrm{good}}$ is a probability expectation and $|\Theta F|=|F|$ pointwise,
\[
\|F\|_{\mathrm{OS}}^2=\langle \Theta F\cdot F\rangle^{\mathrm{good}}
\le \langle |F|^2\rangle^{\mathrm{good}}
\le \|F\|_\infty^2.
\]
Take square roots.
\end{proof}

On the patched-chart state, configurations in each unit block satisfy a fixed exponential chart bound, hence belong to a compact set.
Local functionals are smooth functions on a finite-dimensional compact manifold.

\begin{proposition}[Supremum control by the local seminorm (proved in Proposition~\ref{ass:seminorm_axioms})]\label{ass:sup_from_local}
There exists a constant $C_\infty$ independent of $j$ and $a$ such that for any local functional $F$ supported in a single $j$-block,
\[
\|F\|_{\infty}\le C_\infty\,\|F\|_j.
\]
More generally, for $F$ supported in a $j$-polymer $X$,
\[
\|F\|_\infty\le C_\infty\,\Gamma_j(X)\,\|F\|_j,
\]
where $\Gamma_j(X)$ is the polymer weight used in the definition of $\|\cdot\|_j$.
\end{proposition}

now justified by compactness of the patched-chart configuration space.

Given an insertion activity $J$, define the associated observable
\[
F_J(f):=\sum_{X\in\mathscr P_j} J(X;f),
\]
with the sum restricted to polymers intersecting the $j$-neighborhood of $\mathrm{supp}(f)$ (finite by construction).

\begin{proposition}[Marked-polymer counting (proved in Proposition~\ref{ass:seminorm_axioms})]\label{ass:counting}
There exists $C_{\mathrm{count}}$ independent of $j$ and $a$ such that for any polymer weights $\Gamma_j$,
\[
\sum_{X\ni x} \Gamma_j(X)\ \le\ C_{\mathrm{count}},
\]
uniformly in the lattice site $x$ and uniformly in $j$.
(The same constant controls polymer sums for insertions supported near a fixed compact $K$.)
\end{proposition}

It is purely combinatorial.

\begin{theorem}[Insertion seminorm $\Rightarrow$ OS seminorm]\label{thm:ins_to_OS}\label{lem:ins_to_OS}
Assume \ref{ass:sup_from_local} and \ref{ass:counting}. Then for each dyadic scale $j$ and each insertion activity $J$ supported in the positive-time region,
\[
\|F_J\|_{\mathrm{OS};k,K,\varepsilon}
\le C_{\mathrm{OS}\leftarrow\mathrm{ins}}\ \|J\|_{\mathrm{ins};j;k,K,\varepsilon},
\]
with explicit constant
\begin{equation}\label{eq:C_OS_ins}
C_{\mathrm{OS}\leftarrow\mathrm{ins}}
:= C_\infty\,C_{\mathrm{count}}.
\end{equation}
The constant is independent of $j$ and independent of lattice spacing $a$ under the patched-chart construction.
\end{theorem}

\begin{proof}
Fix $f$ with $\mathrm{supp}(f)\subset K$ and $|f|_{C^k}\le 1$.
By Lemma~\ref{lem:OS_by_sup},
\[
\|F_J(f)\|_{\mathrm{OS}}\le \|F_J(f)\|_\infty \le \sum_X \|J(X;f)\|_\infty.
\]
Apply Proposition~\ref{ass:sup_from_local}:
\[
\|J(X;f)\|_\infty \le C_\infty\,\Gamma_j(X)\,\|J(X;f)\|_j.
\]
By the definition \eqref{eq:ins_norm},
$\|J(X;f)\|_j\le \Gamma_j(X)\,\|J\|_{\mathrm{ins};j;k,K,\varepsilon}$.
Hence
\[
\|F_J(f)\|_{\mathrm{OS}}
\le C_\infty\,\|J\|_{\mathrm{ins};j;k,K,\varepsilon}\ \sum_X \Gamma_j(X).
\]
The sum is over polymers intersecting the fixed neighborhood of $K$; by Proposition~\ref{ass:counting} it is bounded by $C_{\mathrm{count}}$.
Taking the supremum over $f$ yields the claim with constant $C_\infty C_{\mathrm{count}}$.
\end{proof}



\begin{corinsword}[Fixed bounded-word variant of Theorem~\ref{thm:ins_to_OS}]\label{cor:ins_to_OS_bounded_word}
Fix a dyadic scale $j$, a compact $K_f\subset\{x_0\ge \varepsilon\}$, and an insertion activity
$J=\{J(X;f)\}_{X\in\mathscr P_j}$ supported on polymers intersecting the $j$-neighborhood of $K_f$.
Let
\[
\mathbf B=(B_0,\dots,B_q)
\]
be a fixed finite family of bounded positive-time cylinder / flowed observables, and let
$K_{B,f}\subset\{x_0\ge \varepsilon\}$ be a compact neighborhood containing $K_f$ together with the supports of all $B_r$.
For each scale $j$, choose a fixed finite $j$-polymer neighborhood $X_{B,f}(j)$ covering the additional support carried by the word $\mathbf B$ and define the bounded-word insertion activity by
\[
J^{\mathbf B}(X\cup X_{B,f}(j);f):=B_0\,J(X;f)\,B_1\cdots B_q.
\]
Then there exists a constant
\[
C_{\mathrm{word}}(K_{B,f};\mathbf B^{\sharp})\ge 1,
\qquad
\mathbf B^{\sharp}:=(B_0^{\sharp},\dots,B_q^{\sharp}),
\]
depending only on the normalized fixed word $\mathbf B^{\sharp}$ and on the support neighborhood $K_{B,f}$, but independent of $J$, $j$, and $a$, such that
\begin{equation}\label{eq:ins_word_bound}
\|J^{\mathbf B}\|_{\mathrm{ins};j;k,K_{B,f},\varepsilon}
\le
C_{\mathrm{supp}}(K_{B,f})\,
C_{\mathrm{word}}(K_{B,f};\mathbf B^{\sharp})\,
\Bigl(\prod_{r=0}^{q}\|B_r\|_\infty\Bigr)\,
\|J\|_{\mathrm{ins};j;k,K_f,\varepsilon}.
\end{equation}
Consequently,
\begin{equation}\label{eq:OS_word_bound}
\|F_{J^{\mathbf B}}\|_{\mathrm{OS};k,K_{B,f},\varepsilon}
\le
C_{\mathrm{OS}\leftarrow\mathrm{ins}}\,
C_{\mathrm{supp}}(K_{B,f})\,
C_{\mathrm{word}}(K_{B,f};\mathbf B^{\sharp})\,
\Bigl(\prod_{r=0}^{q}\|B_r\|_\infty\Bigr)\,
\|J\|_{\mathrm{ins};j;k,K_f,\varepsilon}.
\end{equation}
\end{corinsword}

\begin{proof}
If one of the factors $B_r$ vanishes then both claims are immediate, so assume $B_r\neq 0$ and write
\[
B_r=\|B_r\|_\infty\,B_r^{\sharp},
\qquad
\|B_r^{\sharp}\|_\infty=1.
\]
Adjoining the fixed support carried by the bounded word enlarges every marked polymer only by the fixed finite neighborhood $X_{B,f}(j)$.
Therefore the associated polymer weight changes by at most a uniform support factor,
\[
\Gamma_j\bigl(X\cup X_{B,f}(j)\bigr)
\le
C_{\mathrm{supp}}(K_{B,f})\,\Gamma_j(X),
\]
with $C_{\mathrm{supp}}(K_{B,f})$ depending only on the chosen support neighborhood and not on $J$, $j$, or $a$.
On the patched-chart compact set, the fixed normalized factors $B_r^{\sharp}$ have bounded local seminorms on $K_{B,f}$; iterating the local product / regularity axioms therefore yields a constant
$C_{\mathrm{word}}(K_{B,f};\mathbf B^{\sharp})$ such that for every marked polymer term,
\[
\bigl\|B_0\,J(X;f)\,B_1\cdots B_q\bigr\|_j
\le
C_{\mathrm{word}}(K_{B,f};\mathbf B^{\sharp})\,
\Bigl(\prod_{r=0}^{q}\|B_r\|_\infty\Bigr)
\|J(X;f)\|_j.
\]
Combining the last two displays and taking the supremum in the insertion norm \eqref{eq:ins_norm} gives \eqref{eq:ins_word_bound}.
Applying Theorem~\ref{thm:ins_to_OS} to the modified activity $J^{\mathbf B}$ then gives \eqref{eq:OS_word_bound}.
\end{proof}

\begin{corollary}[OS control of the insertion RG remainder]\label{cor:A1_bridge}
Let $\mathcal R_{*,t}$ be the insertion RG remainder at matching scale $\mu$ from Theorem~\ref{thm:ins_decomp}.
Under the patched-chart state and the hypotheses of Theorem~\ref{thm:ins_to_OS},
\[
\|\mathcal R_{*,t}\|_{\mathrm{OS};k,K,\varepsilon}
\le C_{\mathrm{OS}\leftarrow\mathrm{ins}}\ \|\mathcal R_{*,t}\|_{\mathrm{ins};\mu;k,K,\varepsilon}.
\]
Thus the insertion-norm remainder bound converts directly into the A1 OS-seminorm remainder bound.
\end{corollary}

\begin{remark}[End of the insertion-to-OS transfer]
Theorem~\ref{thm:ins_to_OS}, Corollary~\ref{cor:ins_to_OS_bounded_word}, and Corollary~\ref{cor:A1_bridge} give an explicit comparison between the insertion seminorm
and the OS seminorm for positive-time insertions supported in $K$, both in the basic single-insertion channel and after adjoining a fixed bounded positive-time word.  In particular, insertion-RG remainder bounds
translate directly into OS-seminorm remainder bounds.
Combined with the coefficient invertibility bounds of Section~4.7 (Corollary~\ref{cor:Lambda}), this completes the
Lane-A small-flow-time convergence argument.  De-patching to the physical Wilson state is handled by
Lemmas~\ref{lem:depatch_block} and \ref{lem:depatch_continuum} in Section~\ref{sec:patching}.
\end{remark}




\paragraph{Packet 5--7 certification note.}
Appendix~\ref{app:normalization_crosscheck}, \S\ref{subsec:LaneA_packet_interface_map}, freezes the Lane~A packet interface in Table~\ref{tab:LaneA_packet_interface_map}.
For the theorem-facing spine used later by Appendix~\ref{app:route1} and Appendix~\ref{app:axioms_package}, that table records three role boundaries that should not be blurred in prose: the present coefficient block exports only the finite-truncation inverse / SFTE package, Theorem~\ref{thm:finite_cap_flowed_sector_bridge} is the first mixed-correlator closure node, and the inductive-union extension plus bounded-base cyclicity appear only at Corollary~\ref{cor:LaneA_full_scope} and Theorem~\ref{thm:LaneA_bounded_base_cyclic}.

\subsection{Auxiliary: coefficient matrix invertibility (diagonal blocks + finite triangular assembly)}
Before entering the coefficient file, it is helpful to separate three roles that are easy to conflate on a first read.
First, the scalar/tensor/loop sectors are retained as \emph{audited sample certificates}: they show concretely how the abstract inverse-control mechanism looks in familiar channels.
For the first dim-$6$ loop block, the paper now writes the audited UV matrix explicitly,
\[
D_{6,\mathrm{tree}}^{PRC}=
\begin{pmatrix}
0&1/8\\[0.15em]
1/24&0
\end{pmatrix},
\]
rather than leaving the determinant certificate in symbolic form.
Second, the load-bearing widening step is the \emph{canonical exact-dimension block formalism}:
for each engineering dimension / CP parity / Euclidean tensor type, the local block is first formalized as an exact-dimension quotient space,
Appendix~\ref{app:normalization_crosscheck} proves that the small-flow substitution induces the identity on that quotient before any basis is chosen,
and only afterwards is the block written in a fixed complete-contraction basis and renormalized coherently across engineering caps.
The auxiliary insertion-RG subsection now proves a \emph{basis-free one-shell transport theorem} for those canonical blocks,
so the later finite-truncation theorem no longer appeals to a placeholder assumption about ``each fixed finite operator basis''.
Third, those canonical blocks are assembled into finite engineering caps and then passed to the inductive union.
Appendix~\ref{app:normalization_crosscheck} now records that basis/normalization protocol explicitly, so a referee can read the full-scope Lane~A statement without reconstructing the bookkeeping from earlier scalar examples.
To pass from flowed observables to sharp renormalized local fields, we use the small-flow-time expansion sectorwise
and invert the resulting finite-dimensional mixing system.
Concretely, in each symmetry sector the flowed basis is related to a renormalized local operator basis by a coefficient matrix
$\mathbf C(t,\mu)$.
We require that for sufficiently small $t>0$ this matrix be invertible with a quantitative bound compatible with the SFTE remainder estimates:
\[
\|\mathbf C(t,\mu)^{-1}\|\le C_{\mathrm{inv}}\,\Lambda(t),
\qquad
\Lambda(t)\,t^\eta\to 0\quad (t\downarrow 0).
\]

\paragraph{How to read the widened Lane~A coefficient package.}
There are three logically distinct layers in this section.
\begin{enumerate}[label=(\roman*),leftmargin=2.3em]
\item \emph{Audited sample sectors.}
We begin with explicit scalar, loop, tensor, and concrete dim-$6$ sample blocks.
These are referee-facing calculations showing how the abstract inverse-control mechanism appears in familiar channels.
\item \emph{Canonical exact-dimension blocks.}
For each engineering dimension, CP parity, and Euclidean tensor type, the local block is written in a complete-contraction basis modulo the quotient relations already built into Definition~\ref{def:local_operator_spaces};
the flowed block is obtained by applying the \emph{same contraction formulas} to $G_{\mu\nu}(t,\cdot)$ and flowed covariant derivatives.
Appendix~\ref{app:normalization_crosscheck} then formalizes these blocks as exact-dimension quotient spaces and proves, \emph{before any basis is chosen},
that the small-flow substitution induces the identity on the associated-graded quotient (Theorem~\ref{thm:LaneA_associated_graded_identity}).
The coherent renormalization of Theorem~\ref{thm:field_algebra_truncation} preserves that intrinsic exact-dimension symbol, and only afterwards do we write matrices in the canonical basis.
\item \emph{Finite truncations and full scope.}
Once each exact-dimension diagonal block has quantitative inverse control, the full finite engineering cap follows from lower-triangular assembly,
and Lane~A then reaches the full sharp-local algebra by the inductive-union argument in Section~\ref{sec:sharp_local}.
\end{enumerate}
Appendix~\ref{app:normalization_crosscheck} now records this basis/normalization protocol explicitly.
The explicit scalar/tensor/loop computations therefore remain as normalization cross-checks and audited examples,
while the load-bearing widening step is the canonical finite-truncation theorem near the end of the section.

\medskip
\begin{itemize}
\item CP-even scalar dim-4 (after centering): coefficient $c_S(t,\mu)$ in $\mathcal O_{S,t}\mapsto \mathcal O_S^{\mathrm{ren}}$,
\item CP-odd pseudoscalar dim-4: coefficient $c_Q(t,\mu)$ in $\mathcal O_{Q,t}\mapsto \mathcal O_Q^{\mathrm{ren}}$,
\item small (symmetrized) Wilson loops: coefficient $d_S(t,\mu)$ in $(W_t-1)/\ell^4\mapsto \mathcal O_S^{\mathrm{ren}}$,
\item traceless-symmetric dim-4 tensor probes built from flowed curvature bilinears: finite block $C_T(t,\mu)$ in $\mathbf U_{T,t}\mapsto \mathbf T^{\mathrm{ren}}$,
\item a concrete plaquette/rectangle/chair CP-even dim-$6$ scalar loop block: finite block $D_6^{PRC}(t,\mu)$ in $\mathbf V_{6,t}^{PRC}\mapsto \mathbf O_{6}^{PRC,\mathrm{ren}}$.
\end{itemize}
In each explicitly stabilized case, the inverse is bounded by a \textbf{polylogarithmic} function $\Lambda(t)$,
so $\Lambda(t)t^\eta\to 0$ for every $\eta>0$.
For the concrete dim-$6$ loop block below, we record the finite-dimensional UV/tree-level certificate explicitly:
\[
\begin{aligned}
D_{6,\mathrm{tree}}^{PRC}
&=
\begin{pmatrix}
0 & 1/8\\[0.15em]
1/24 & 0
\end{pmatrix},\\
\det D_{6,\mathrm{tree}}^{PRC}
&=-\frac{1}{192}.
\end{aligned}
\]
and the final theorems in this section then show how the same inverse-control mechanism extends from these audited examples to arbitrary fixed symmetry blocks and hence to every finite sharp-local truncation.

\paragraph{Standing convention (physical Wilson state; chart truncation as an analytic device).}
All expectations are taken in the physical Wilson lattice state.  We occasionally pass through a good-flow chart truncation
to obtain uniform kernel and RG bounds; the truncation error is removed by the de-patching stability lemmas of
Section~\ref{sec:patching} (Lemmas~\ref{lem:depatch_block} and \ref{lem:depatch_continuum}).

In the three target sectors, after the standard subtractions, the mixing matrix is $1\times 1$.

The flowed probe is $\mathcal O_{S,t}=\tfrac14\mathrm{tr}\,G_{\mu\nu}(t)G_{\mu\nu}(t)$.
The only lower-dimension mixing is with the identity; centering removes it:
\[
\mathcal O_{S,t}(f)-\langle \mathcal O_{S,t}(f)\rangle\mathbf 1
= c_S(t,\mu)\,\mathcal O_S^{\mathrm{ren}}(f;\mu) + R_{S,t}^{\circ}(f).
\]

\[
\mathcal O_{Q,t}(f)=c_Q(t,\mu)\,\mathcal O_Q^{\mathrm{ren}}(f;\mu) + R_{Q,t}(f).
\]

For a symmetrized family of flowed plaquettes of side $\ell$ (averaged over orientations) define
\[
\mathcal W_t^{(\ell)}(x):=\frac{1}{6}\sum_{\mu<\nu}\frac{1}{\dim(\rho)}\Re\Tr_{\rho}\, \mathcal P \exp\Big(\oint_{C^{\mu\nu}_\ell} B(t,\cdot)\cdot d x\Big),
\qquad
\ell^2=\kappa_W t.
\]
Then
\[
\frac{\mathcal W_t^{(\ell)}(f)-\langle \mathcal W_t^{(\ell)}(f)\rangle}{\ell^4}
=
d_S(t,\mu)\,\mathcal O_S^{\mathrm{ren}}(f;\mu) + R_{W,t}(f).
\]
Thus in each sector, $\mathbf C(t,\mu)$ reduces to a scalar coefficient and invertibility is simply nonvanishing of that coefficient.

We use the insertion RG machinery from Theorem~\ref{thm:ins_decomp} and Definition~\ref{def:RG_ins}. Fix a renormalization scale $\mu$ in the AF window (so $u(\mu)\le u_\ast$),
and define the ``matching'' step number
\[
n(t,\mu)=\Big\lceil \log_2\Big(\frac{1}{\mu\sqrt t}\Big)\Big\rceil.
\]
The coefficients $c_S(t,\mu)$, $c_Q(t,\mu)$, $d_S(t,\mu)$ are defined as the extracted relevant coefficients at the matching scale
in the insertion RG decomposition ( Theorem~\ref{thm:ins_decomp}).

For each sector, define the one-step multiplier $M_{*,j}$ as follows:
start with the normalized basis insertion $\Phi_{*,j}$ at scale $j$ representing $\mathcal O_*^{\mathrm{ren}}$ in that sector;
apply one good-event (chart-truncated) insertion RG step ( Definition~\ref{def:RG_ins});
extract the coefficient of $\Phi_{*,j+1}$ in the relevant sector at scale $j+1$.
This yields
\[
\Phi_{*,j}\ \xrightarrow{\mathcal R^{\mathrm{ins,good}}_j}\ M_{*,j}\,\Phi_{*,j+1}\ +\ (\text{extracted-irrelevant}).
\]
The number $M_{*,j}$ depends on $(u_j,K_j)$ but is \emph{finite-dimensional}.

\begin{proposition}[One-step multiplier bound]\label{prop:mult_bound}
There exist constants $u_\ast>0$ and $C_{M,S},C_{M,Q},C_{M,W}$ such that whenever the bulk RG lies in the AF ball
($0\le u_j\le u_\ast$ and $\|K_j\|_j\le r_K$),
\begin{equation}\label{eq:M_bounds}
|M_{S,j}-1|\le C_{M,S}\,u_j,\qquad
|M_{Q,j}-1|\le C_{M,Q}\,u_j,\qquad
|M_{W,j}-1|\le C_{M,W}\,u_j.
\end{equation}
We additionally require $u_\ast\le \min\{(2C_{M,S})^{-1},(2C_{M,W})^{-1}\}$ so that $|M_{S,j}|,|M_{W,j}|\ge \tfrac12$.
\end{proposition}

\begin{proof}
A proof (with constants tracked in terms of the insertion integration bound, extraction/projection bounds, and the stable-manifold forcing constant)
is given in Appendix~\ref{app:multiplier_bounds}.
\end{proof}

\begin{remark}[Optional strengthening in the CP-odd channel]
If one inputs additional information about the vanishing of the leading anomalous dimension in the CP-odd channel,
the middle bound in \eqref{eq:M_bounds} can be improved to $|M_{Q,j}-1|=O(u_j^2)$.
The present paper does not use this improvement.
\end{remark}


\paragraph{Comment (status).}
These are finite-dimensional consequences of analyticity of the insertion RG map on the AF ball (as inherited from Theorem~\ref{thm:RG_analytic_bounds}) and normalization $M(0,0)=1$.
Let $u_j$ denote the dyadic running coupling at scale $j$ in the AF window.
\begin{equation}\label{eq:beta}
u_{j+1}=u_j + A\,u_j^2 + r_j,\qquad A:=2b_0\log 2>0,\qquad |r_j|\le B\,u_j^3,
\end{equation}
with $B$ scale-independent.

\begin{lemma}[Sum of couplings grows like $\log(u_{\mathrm{IR}}/u_{\mathrm{UV}})$]\label{lem:sum_u}
Fix $u_\ast>0$ such that $B u_\ast\le A/2$ and assume $u_j\le u_\ast$ for $j=j_0,\dots,j_0+N$.
Then
\begin{equation}\label{eq:sum_u}
\sum_{m=0}^{N-1} u_{j_0+m}
\ \le\ \frac{2}{A_-}\ \log\Big(\frac{u_{j_0+N}}{u_{j_0}}\Big),
\qquad
A_-:=\frac{A}{1+A u_\ast}.
\end{equation}
In particular, if $u_{j_0+N}=u(\mu)$ is fixed and $u_{j_0}=u(\mu_t)\to 0$ as $t\downarrow 0$, then the sum diverges at most like $\log(1/u(\mu_t))$.
\end{lemma}

\begin{proof}
From \eqref{eq:beta} and $B u_\ast\le A/2$, we have $u_{j+1}\ge u_j + (A/2)u_j^2$.
Thus
\[
\frac{1}{u_{j+1}}=\frac{1}{u_j(1+(A/2)u_j)}\le \frac{1}{u_j}-\frac{A/2}{1+(A/2)u_\ast}\le \frac{1}{u_j}-\frac{A_-}{2}.
\]
Iterating yields $u_{j_0+m}\le (u_{j_0}^{-1}-\tfrac{A_-}{2}m)^{-1}$. Summing the harmonic bound gives \eqref{eq:sum_u}.
\end{proof}

Define the RG coefficient at matching scale $\mu$ by
\[
c_S(t,\mu):=\prod_{m=0}^{n(t,\mu)-1} M_{S,j_t+m}\cdot c_{S,\mathrm{UV}}(t),
\]
where $c_{S,\mathrm{UV}}(t)$ is the ``initial'' coefficient at the flow-matching UV scale $j_t$ (by normalization $c_{S,\mathrm{UV}}(t)=1+O(u(\mu_t))$).

\begin{theorem}[CP-even nondegeneracy and polylog inverse]\label{thm:cS}
Assume Proposition~\ref{prop:mult_bound} and that the RG trajectory from scale $\mu_t\simeq 1/\sqrt t$ down to $\mu$ stays in the AF ball ($u_j\le u_\ast$).
Then for all sufficiently small $t$,
\begin{equation}\label{eq:cS_bounds}
0< c_{S,\min}(t,\mu)\ \le\ |c_S(t,\mu)|\ \le\ c_{S,\max}(t,\mu),
\end{equation}
with explicit polylog bounds
\begin{equation}\label{eq:polylog}
|c_S(t,\mu)|^{-1}\ \le\ C_{S,\mathrm{inv}}\ \Big(\frac{u(\mu)}{u(\mu_t)}\Big)^{p_S},
\qquad
p_S:=\frac{4C_{M,S}}{A_-},
\end{equation}
and $C_{S,\mathrm{inv}}$ depending only on the UV normalization bound $|c_{S,\mathrm{UV}}(t)|\ge 1/2$ for small $t$.
In particular, since $u(\mu_t)\downarrow 0$ as $t\downarrow 0$ in an asymptotically free trajectory,
\[
\Lambda_S(t):=|c_S(t,\mu)|^{-1}
\]
grows at most polylogarithmically in $t^{-1}$.
\end{theorem}

\begin{proof}
For $u_j\le u_\ast$ with $u_\ast\le (2C_{M,S})^{-1}$, we have $|M_{S,j}|\ge 1/2$ and
$|\log|M_{S,j}||\le 2|M_{S,j}-1|\le 2C_{M,S}u_j$.
Hence
\[
|\log|c_S(t,\mu)|-\log|c_{S,\mathrm{UV}}(t)||\le 2C_{M,S}\sum_{m=0}^{n-1}u_{j_t+m}.
\]
Use Lemma~\ref{lem:sum_u} and exponentiate, absorbing $|c_{S,\mathrm{UV}}(t)|^{-1}\le 2$ into $C_{S,\mathrm{inv}}$.
\end{proof}

\begin{theorem}[CP-odd nondegeneracy and polylog inverse]\label{thm:cQ}
Assume Proposition~\ref{prop:mult_bound} and that the RG trajectory from scale $\mu_t\simeq 1/\sqrt t$ down to $\mu$ stays in the AF ball ($u_j\le u_\ast$).
Then for all sufficiently small $t$,
\begin{equation}\label{eq:cQ_bounds}
0< c_{Q,\min}(t,\mu)\ \le\ |c_Q(t,\mu)|\ \le\ c_{Q,\max}(t,\mu),
\end{equation}
with explicit polylog bounds
\begin{equation}\label{eq:polylog_Q}
|c_Q(t,\mu)|^{-1}\ \le\ C_{Q,\mathrm{inv}}\ \Big(\frac{u(\mu)}{u(\mu_t)}\Big)^{p_Q},
\qquad
p_Q:=\frac{4C_{M,Q}}{A_-},
\end{equation}
and $C_{Q,\mathrm{inv}}$ depending only on the UV normalization bound $|c_{Q,\mathrm{UV}}(t)|\ge 1/2$ for small $t$.
In particular, $\Lambda_Q(t):=|c_Q(t,\mu)|^{-1}$ grows at most polylogarithmically in $t^{-1}$, hence $\Lambda_Q(t)t^\eta\to0$ for every $\eta>0$.
\end{theorem}

\begin{proof}
From \eqref{eq:M_bounds} we have $|M_{Q,j}-1|\le C_{M,Q}u_j$.
For $u_j\le u_\ast\le (2C_{M,Q})^{-1}$ we have $|\log|M_{Q,j}||\le 2C_{M,Q}u_j$.
Therefore
\[
\Big|\log\Big|\prod_{m=0}^{n-1} M_{Q,j_t+m}\Big|\Big|
\le 2C_{M,Q}\sum_{m=0}^{n-1} u_{j_t+m}.
\]
Apply Lemma~\ref{lem:sum_u} to bound the sum by $\frac{4}{A_-}\log\!\big(\frac{u(\mu)}{u(\mu_t)}\big)$, and absorb the UV normalization factor $c_{Q,\mathrm{UV}}(t)=1+O(u(\mu_t))$ into $C_{Q,\mathrm{inv}}$ for $t$ small.
\end{proof}


At tree level (classical expansion), the symmetrized plaquette probe satisfies
\begin{equation}\label{eq:dS0}
\frac{\mathcal W^{(\ell)}_0(x)-1}{\ell^4}
=\frac{1}{6N}\,\mathcal O_S(x)\ +\ O(\ell^2),
\end{equation}
so the leading coefficient is
\[
d_S^{(0)}=\frac{1}{6N}\neq 0.
\]
Quantum and flow corrections renormalize this by the same RG running as $\mathcal O_S$ plus $O(u)$ perturbations.

\begin{theorem}[Loop coefficient nondegeneracy and inverse bound]\label{thm:dS}
Assume Proposition~\ref{prop:mult_bound} and the AF-ball condition along the trajectory.
Then for sufficiently small $t$,
\[
|d_S(t,\mu)|\ \ge\ \frac{|d_S^{(0)}|}{4}\ \Big(\frac{u(\mu_t)}{u(\mu)}\Big)^{p_W},
\qquad
|d_S(t,\mu)|^{-1}\ \le\ C_{W,\mathrm{inv}}\ \Big(\frac{u(\mu)}{u(\mu_t)}\Big)^{p_W},
\]
with exponent $p_W:=\frac{4C_{M,W}}{A_-}$ and $C_{W,\mathrm{inv}}$ depending only on the UV normalization and the choice of loop family.
In particular, $|d_S(t,\mu)|^{-1}$ grows at most polylogarithmically as $t\downarrow 0$.
\end{theorem}

\begin{proof}
Normalize the UV loop coefficient so that $d_{S,\mathrm{UV}}(t)$ lies within a factor $2$ of $d_S^{(0)}$ for $t$ sufficiently small
(using the classical expansion \eqref{eq:dS0} and $u(\mu_t)\to 0$).
Then run the same multiplier-product bound as in Theorem~\ref{thm:cS} with $M_{W,j}$.
\end{proof}

\begin{corollary}[Sectorwise invertibility condition]\label{cor:Lambda}
In each of the three sectors, set $\Lambda(t)$ to be the corresponding inverse bound:
\[
\Lambda_S(t)=|c_S(t,\mu)|^{-1},\qquad
\Lambda_Q(t)=|c_Q(t,\mu)|^{-1},\qquad
\Lambda_W(t)=|d_S(t,\mu)|^{-1}.
\]
Then $\Lambda(t)$ grows at most polylogarithmically in $t^{-1}$, hence for every $\eta>0$,
\[
\Lambda(t)\,t^\eta\ \xrightarrow[t\downarrow 0]{}\ 0.
\]
Therefore the required small-flow-time invertibility criterion is satisfied sectorwise for the CP-even, CP-odd, and loop generators.
\end{corollary}

\begin{proof}
$t^\eta (\log(1/t))^p\to 0$ as $t\downarrow0$ for all fixed $\eta>0$ and $p\ge 0$.
\end{proof}



\paragraph{A generic finite-dimensional diagonal-block criterion.}
The scalar multiplier-product proofs above use only two ingredients:
near-identity one-step transfer in the AF ball and an invertible UV matching coefficient.
The same mechanism closes any fixed finite symmetry block.
For the exact-dimension quotient blocks attached to Definition~\ref{def:local_operator_spaces}, the needed coefficient extraction is now theorem-level rather than schematic:
Appendix~\ref{app:normalization_crosscheck} proves that quotienting by lower-dimension terms and homogeneous relations commutes with the small-flow substitution and yields a literal identity matrix on every arbitrary exact-dimension block.
Thus the proposition below is applied to canonically defined finite-dimensional blocks, not to a sector-by-sector template.

\begin{proposition}[Finite-dimensional one-step block multiplier bound]\label{prop:block_mult_bound}
Let $\mathfrak s$ be a fixed finite symmetry sector with basis size $m_{\mathfrak s}<\infty$
(for instance a traceless-symmetric tensor channel or the concrete plaquette/rectangle/chair dim-$6$ scalar block treated below).
Let
\[
\mathbf\Phi_{\mathfrak s,j}
=
\bigl(\Phi_{\mathfrak s,j}^{(1)},\dots,\Phi_{\mathfrak s,j}^{(m_{\mathfrak s})}\bigr)^{\mathsf T}
\]
be the normalized basis insertion vector at scale $j$ and define the extracted one-step transfer matrix
$M_{\mathfrak s,j}\in \mathrm{Mat}_{m_{\mathfrak s}}(\mathbb R)$ by
\[
\mathbf\Phi_{\mathfrak s,j}
\xrightarrow{\mathcal R^{\mathrm{ins,good}}_j}
M_{\mathfrak s,j}\,\mathbf\Phi_{\mathfrak s,j+1}
+ (\text{extracted-irrelevant}).
\]
Then there exists a constant $C_{M,\mathfrak s}<\infty$ such that whenever the bulk RG lies in the AF ball
($0\le u_j\le u_\ast$ and $\|K_j\|_j\le r_K$),
\begin{equation}\label{eq:block_mult_near_identity}
\|M_{\mathfrak s,j}-I\|_{\mathrm{op}}\le C_{M,\mathfrak s}\,u_j.
\end{equation}
After shrinking $u_\ast$ so that $C_{M,\mathfrak s}u_\ast\le \tfrac12$, each $M_{\mathfrak s,j}$ is invertible and
\begin{equation}\label{eq:block_mult_inverse_bound}
\|M_{\mathfrak s,j}^{-1}\|_{\mathrm{op}}
\le (1-C_{M,\mathfrak s}u_j)^{-1}
\le 1+2C_{M,\mathfrak s}u_j
\le \exp(2C_{M,\mathfrak s}u_j).
\end{equation}
\end{proposition}

\begin{proof}
For each input basis vector $\Phi_{\mathfrak s,j}^{(q)}$ and each output coefficient index $p$,
let $\mathsf P^{(p)}_{\mathfrak s,j+1}$ denote the bounded projection extracting the coefficient of
$\Phi_{\mathfrak s,j+1}^{(p)}$ after one insertion RG step.
Exactly as in Appendix~\ref{app:multiplier_bounds}, Lemma~\ref{lem:ins_step_Lip} gives
\[
\big\|\mathcal R^{\mathrm{ins,good}}_j(u,K,\Phi_{\mathfrak s,j}^{(q)})
-\mathcal R^{\mathrm{ins,good}}_j(0,0,\Phi_{\mathfrak s,j}^{(q)})\big\|_{j+1}
\le C_{\mathrm{Lip},\mathfrak s}\,\bigl(u+\|K\|_j\bigr),
\]
with a constant uniform in $j$ because the sector dimension is fixed.
Hence
\[
\big|(M_{\mathfrak s,j}-I)_{pq}\big|
\le C_{\mathrm{proj},\mathfrak s}\,C_{\mathrm{Lip},\mathfrak s}\,\bigl(u_j+\|K_j\|_j\bigr).
\]
On the AF ball the stable-manifold forcing bound gives
$\|K_j\|_j\le \frac{C_F}{1-\kappa}\,u_\ast\,u_j$.
Because $m_{\mathfrak s}<\infty$, operator norm and max-entry norm are equivalent, so after absorbing the fixed sector-size factor we obtain
\eqref{eq:block_mult_near_identity}.
If $\|M_{\mathfrak s,j}-I\|_{\mathrm{op}}\le \tfrac12$, the Neumann-series bound gives
$\|M_{\mathfrak s,j}^{-1}\|_{\mathrm{op}}\le (1-\|M_{\mathfrak s,j}-I\|_{\mathrm{op}})^{-1}$,
which implies \eqref{eq:block_mult_inverse_bound}.
\end{proof}

\begin{theorem}[Generic finite-dimensional diagonal-block nondegeneracy]\label{thm:block_polylog}
Let $\mathfrak s$ be a fixed finite symmetry block and let
\[
C_{\mathfrak s}(t,\mu)\in \mathrm{Mat}_{m_{\mathfrak s}}(\mathbb R)
\]
be the coefficient matrix in the small-flow-time expansion from flowed probes to renormalized local operators in that block.
Assume:
\begin{enumerate}[label=(\roman*),leftmargin=2.3em]
\item the RG trajectory from the matching scale $\mu_t\simeq 1/\sqrt t$ down to $\mu$ stays in the AF ball;
\item the one-step transfer matrices satisfy Proposition~\ref{prop:block_mult_bound};
\item the UV matching matrix has a nondegenerate limit,
\begin{equation}\label{eq:block_UV_nondeg}
C_{\mathfrak s,\mathrm{UV}}(t)=C_{\mathfrak s}^{(0)}+O\!\big(u(\mu_t)\big),
\qquad
C_{\mathfrak s}^{(0)}\in \mathrm{GL}_{m_{\mathfrak s}}(\mathbb R).
\end{equation}
\end{enumerate}
Then for all sufficiently small $t$ the block $C_{\mathfrak s}(t,\mu)$ is invertible and obeys the polylogarithmic inverse bound
\begin{equation}\label{eq:block_polylog_bound}
\|C_{\mathfrak s}(t,\mu)^{-1}\|_{\mathrm{op}}
\le
C_{\mathfrak s,\mathrm{inv}}
\Big(\frac{u(\mu)}{u(\mu_t)}\Big)^{p_{\mathfrak s}},
\qquad
p_{\mathfrak s}:=\frac{4C_{M,\mathfrak s}}{A_-},
\end{equation}
for a constant $C_{\mathfrak s,\mathrm{inv}}<\infty$ depending only on
$\|(C_{\mathfrak s}^{(0)})^{-1}\|_{\mathrm{op}}$ and the small-$t$ UV normalization.
In particular, $\|C_{\mathfrak s}(t,\mu)^{-1}\|_{\mathrm{op}}$ grows at most polylogarithmically as $t\downarrow 0$.
\end{theorem}

\begin{proof}
Write $n=n(t,\mu)$ and factor the coefficient block as
\[
C_{\mathfrak s}(t,\mu)
=
M_{\mathfrak s,j_t+n-1}\cdots M_{\mathfrak s,j_t}\,C_{\mathfrak s,\mathrm{UV}}(t).
\]
By \eqref{eq:block_UV_nondeg}, $C_{\mathfrak s,\mathrm{UV}}(t)$ is invertible for all sufficiently small $t$ and
$\|C_{\mathfrak s,\mathrm{UV}}(t)^{-1}\|_{\mathrm{op}}\le 2\|(C_{\mathfrak s}^{(0)})^{-1}\|_{\mathrm{op}}$ after shrinking $t_0$ if necessary.
Therefore
\[
\|C_{\mathfrak s}(t,\mu)^{-1}\|_{\mathrm{op}}
\le
\|C_{\mathfrak s,\mathrm{UV}}(t)^{-1}\|_{\mathrm{op}}
\prod_{m=0}^{n-1}\|M_{\mathfrak s,j_t+m}^{-1}\|_{\mathrm{op}}.
\]
Using \eqref{eq:block_mult_inverse_bound},
\[
\prod_{m=0}^{n-1}\|M_{\mathfrak s,j_t+m}^{-1}\|_{\mathrm{op}}
\le
\exp\!\Big(2C_{M,\mathfrak s}\sum_{m=0}^{n-1}u_{j_t+m}\Big).
\]
Apply Lemma~\ref{lem:sum_u} to bound the sum by
$\frac{2}{A_-}\log\!\big(\frac{u(\mu)}{u(\mu_t)}\big)$ and absorb the UV prefactor into
$C_{\mathfrak s,\mathrm{inv}}$.
This yields \eqref{eq:block_polylog_bound}.
\end{proof}

\paragraph{New diagonal block: traceless-symmetric dim-4 tensor channel.}
Define the flowed traceless tensor probe
\[
U_{\mu\nu,t}(x)
:=
\tr\!\Big(
G_{\mu\lambda}(t,x)G_{\nu\lambda}(t,x)
-\frac14\delta_{\mu\nu}\,G_{\rho\sigma}(t,x)G_{\rho\sigma}(t,x)
\Big).
\]
Smearing against traceless symmetric test tensors $h^{\mu\nu}$ gives a finite family $\mathbf U_{T,t}(h)$.
Because the channel is traceless symmetric, there is no mixing with the identity or the scalar $F^2$ sector;
the diagonal block is therefore the renormalized traceless-tensor family $\mathbf T^{\mathrm{ren}}(h;\mu)$ itself:
\[
\mathbf U_{T,t}(h)
=
C_T(t,\mu)\,\mathbf T^{\mathrm{ren}}(h;\mu)+\mathbf R_{T,t}(h).
\]
By normalization of the tensor basis at the Gaussian point, the UV matching matrix satisfies
\[
C_{T,\mathrm{UV}}(t)=I_{m_T}+O\!\big(u(\mu_t)\big).
\]

\begin{theorem}[Traceless-tensor diagonal-block inverse control]\label{thm:cT_block}
Assume the AF-ball hypothesis along the trajectory and the one-step block multiplier bound of
Proposition~\ref{prop:block_mult_bound} in the traceless-symmetric dim-4 tensor channel.
Then for all sufficiently small $t$,
\begin{equation}\label{eq:cT_polylog}
\|C_T(t,\mu)^{-1}\|_{\mathrm{op}}
\le
C_{T,\mathrm{inv}}
\Big(\frac{u(\mu)}{u(\mu_t)}\Big)^{p_T},
\qquad
p_T:=\frac{4C_{M,T}}{A_-},
\end{equation}
so $\Lambda_T(t):=\|C_T(t,\mu)^{-1}\|_{\mathrm{op}}$ grows at most polylogarithmically in $t^{-1}$.
In particular, $\Lambda_T(t)\,t^\eta\to 0$ for every $\eta>0$.
\end{theorem}

\begin{proof}
Apply Theorem~\ref{thm:block_polylog} with $\mathfrak s=T$ and $C_T^{(0)}=I_{m_T}$.
\end{proof}

\paragraph{New diagonal block: concrete plaquette/rectangle/chair dim-$6$ scalar loop sector.}
Let $X_{P,t}(f)$, $X_{R,t}(f)$, and $X_{C,t}(f)$ denote centered, orientation-averaged flowed loop densities supported in a fixed $2\ell\times 2\ell$ box, with
\[
\ell^2=\kappa_W t,
\]
corresponding respectively to the plaquette, the planar $1\times 2$ rectangle, and the chair loop.
Normalize the three families so that their common tree-level dim-$4$ scalar coefficient is the same constant $d_S^{(0)}$.
Define the improved CP-even dim-$6$ probes
\[
L_{6,+,t}^{(1)}(f):=\ell^{-2}\bigl(X_{R,t}(f)-X_{P,t}(f)\bigr),
\qquad
L_{6,+,t}^{(2)}(f):=\ell^{-2}\bigl(X_{C,t}(f)-X_{P,t}(f)\bigr).
\]
Because the dim-$4$ contribution cancels by construction, both probes have leading engineering dimension $6$ and live in the CP-even scalar sector.
Write
\[
\mathbf V_{6,t}^{PRC}(f):=
\begin{pmatrix}
L_{6,+,t}^{(1)}(f)\\[0.35em]
L_{6,+,t}^{(2)}(f)
\end{pmatrix}.
\]


\begin{proposition}[Finite UV determinant certificate for the plaquette/rectangle/chair dim-$6$ block]\label{prop:D6_tree_certificate}
For this audited sample block choose the on-shell CP-even dim-$6$ gauge basis
\[
\mathbf O_{6}^{\mathrm{ren}}(f;\mu)
:=
\begin{pmatrix}
\mathcal G_{6}^{(1),\mathrm{ren}}(f;\mu)\\[0.35em]
\mathcal G_{6}^{(2),\mathrm{ren}}(f;\mu)
\end{pmatrix},
\]
where
\[
\mathcal G_{6}^{(1)}\sim \sum_{\mu,\nu,\rho}\Tr\big([D_\mu,F_{\nu\rho}]\,[D_\mu,F_{\nu\rho}]\big),
\qquad
\mathcal G_{6}^{(2)}\sim \sum_{\mu,\nu}\Tr\big([D_\mu,F_{\mu\nu}]\,[D_\mu,F_{\mu\nu}]\big).
\]
This on-shell gauge basis is used only for the explicit certificate; Corollary~\ref{cor:D6_basis_transfer} later transports the same inverse control to any other fixed basis in the same dim-$6$ scalar block.
Then the tree-level matching of the improved probes against that basis has the form
\begin{equation}\label{eq:D6_tree_matrix}
\mathbf V_{6,t}^{PRC}(f)
=
D_{6,\mathrm{tree}}^{PRC}\,\mathbf O_{6}^{\mathrm{ren}}(f;\mu_t)
+\mathbf r_{6,t}^{PRC}(f),
\qquad
\mu_t\simeq 1/\sqrt t,
\end{equation}
with
\[
\|\mathbf r_{6,t}^{PRC}(f)\|=O(\ell^2)+O\!\big(u(\mu_t)\big)
\]
and with the explicit fixed matrix
\begin{equation}\label{eq:D6_tree_det_matrix}
D_{6,\mathrm{tree}}^{PRC}
=
\begin{pmatrix}
0 & \frac18\\[0.45em]
\frac1{24} & 0
\end{pmatrix}.
\end{equation}
Hence the finite UV determinant certificate is
\begin{equation}\label{eq:D6_tree_det}
\Delta_{PRC}:=\det D_{6,\mathrm{tree}}^{PRC}
=
-\frac{1}{192}\neq 0.
\end{equation}
In particular, $D_{6,\mathrm{tree}}^{PRC}\in GL_2(\mathbb R)$.
\end{proposition}

\begin{proof}
Use the standard tree-level Symanzik / L\"uscher--Weisz local matching formulas for the normalized four-link / six-link gauge-loop basis \cite{LuscherWeisz1985,Husung2023}:
for loop coefficients $(e_0,e_1,e_2,e_3)$ satisfying
\[
e_0+8e_1+8e_2+16e_3=1,
\]
the two CP-even dim-$6$ gauge coefficients in the on-shell basis above are
\[
\bar c_1=\frac{e_2}{3},
\qquad
\bar c_2=e_1-e_3+\frac1{12}.
\]
After the common dim-$4$ normalization adopted for $X_{P,t},X_{R,t},X_{C,t}$, the plaquette, rectangle, and chair densities correspond respectively to
\[
P:(e_0,e_1,e_2,e_3)=(1,0,0,0),
\qquad
R:(0,\tfrac18,0,0),
\qquad
C:(0,0,\tfrac18,0).
\]
Hence their tree-level dim-$6$ coefficients are
\[
P:(\bar c_1,\bar c_2)=\Bigl(0,\frac1{12}\Bigr),
\qquad
R:(\bar c_1,\bar c_2)=\Bigl(0,\frac5{24}\Bigr),
\qquad
C:(\bar c_1,\bar c_2)=\Bigl(\frac1{24},\frac1{12}\Bigr).
\]
At matching scale $\mu_t\simeq t^{-1/2}$ this gives
\[
X_{P,t}(f)=d_S^{(0)}\,\mathcal O_S^{\mathrm{ren}}(f;\mu_t)+\ell^2\Bigl(\frac1{12}\,\mathcal G_{6}^{(2),\mathrm{ren}}(f;\mu_t)\Bigr)+r_{P,t}(f),
\]
\[
X_{R,t}(f)=d_S^{(0)}\,\mathcal O_S^{\mathrm{ren}}(f;\mu_t)+\ell^2\Bigl(\frac5{24}\,\mathcal G_{6}^{(2),\mathrm{ren}}(f;\mu_t)\Bigr)+r_{R,t}(f),
\]
\[
X_{C,t}(f)=d_S^{(0)}\,\mathcal O_S^{\mathrm{ren}}(f;\mu_t)+\ell^2\Bigl(\frac1{24}\,\mathcal G_{6}^{(1),\mathrm{ren}}(f;\mu_t)+\frac1{12}\,\mathcal G_{6}^{(2),\mathrm{ren}}(f;\mu_t)\Bigr)+r_{C,t}(f),
\]
with each remainder $r_{*,t}(f)=O(\ell^4)+O(\ell^2u(\mu_t))$.
Subtracting the common plaquette term and dividing by $\ell^2$ yields
\[
L_{6,+,t}^{(1)}(f)=\frac18\,\mathcal G_{6}^{(2),\mathrm{ren}}(f;\mu_t)+O(\ell^2)+O\!\big(u(\mu_t)\big),
\]
\[
L_{6,+,t}^{(2)}(f)=\frac1{24}\,\mathcal G_{6}^{(1),\mathrm{ren}}(f;\mu_t)+O(\ell^2)+O\!\big(u(\mu_t)\big).
\]
This is exactly \eqref{eq:D6_tree_matrix} with matrix \eqref{eq:D6_tree_det_matrix}; taking the determinant gives \eqref{eq:D6_tree_det}.
\end{proof}

Choose the loop-adapted renormalized basis
\[
\mathbf O_{6}^{PRC,\mathrm{ren}}(f;\mu)
:=
D_{6,\mathrm{tree}}^{PRC}\,\mathbf O_{6}^{\mathrm{ren}}(f;\mu).
\]
Because $D_{6,\mathrm{tree}}^{PRC}$ is fixed and invertible, this is just a $t$-independent finite change of basis.
In that basis the UV matching block is normalized to the identity at tree level:
\[
\mathbf V_{6,t}^{PRC}(f)
=
D_6^{PRC}(t,\mu)\,\mathbf O_{6}^{PRC,\mathrm{ren}}(f;\mu)+\mathbf R_{6,t}^{PRC}(f),
\qquad
D_{6,\mathrm{UV}}^{PRC}(t)=I_2+E_{6}^{PRC}(t),
\]
with
\[
\|E_{6}^{PRC}(t)\|_{\mathrm{op}}\le C_{6,\mathrm{UV}}^{PRC}\,u(\mu_t)
\]
for all sufficiently small $t$.

\begin{theorem}[Concrete plaquette/rectangle/chair dim-$6$ block inverse control]\label{thm:D6_block}
Assume:
\begin{enumerate}[label=(\roman*),leftmargin=2.3em]
\item the AF-ball hypothesis holds along the matching trajectory;
\item the one-step block multiplier bound of Proposition~\ref{prop:block_mult_bound} holds in this dim-$6$ scalar block.
\end{enumerate}
Then, after shrinking $t_0$ so that $C_{6,\mathrm{UV}}^{PRC}u(\mu_t)\le \tfrac12$ on $(0,t_0]$, the UV block is uniformly invertible with
\begin{equation}\label{eq:D6_UV_inv}
\|(D_{6,\mathrm{UV}}^{PRC}(t))^{-1}\|_{\mathrm{op}}\le 2.
\end{equation}
Consequently, for all sufficiently small $t$,
\begin{equation}\label{eq:D6_polylog}
\|D_6^{PRC}(t,\mu)^{-1}\|_{\mathrm{op}}
\le
C_{6,\mathrm{inv}}
\Big(\frac{u(\mu)}{u(\mu_t)}\Big)^{p_6},
\qquad
p_6:=\frac{4C_{M,6}}{A_-},
\end{equation}
so the inverse control in this block is again polylogarithmic.
In particular, $\Lambda_6(t):=\|D_6^{PRC}(t,\mu)^{-1}\|_{\mathrm{op}}$ satisfies $\Lambda_6(t)t^\eta\to 0$ for every $\eta>0$.
\end{theorem}

\begin{proof}
The uniform UV inverse bound \eqref{eq:D6_UV_inv} is the Neumann-series estimate for $I_2+E_{6}^{PRC}(t)$ once $\|E_{6}^{PRC}(t)\|_{\mathrm{op}}\le \tfrac12$.
Now apply Theorem~\ref{thm:block_polylog} with $\mathfrak s=6$ and with UV normalization constant coming from \eqref{eq:D6_UV_inv}.
\end{proof}

\begin{corollary}[Basis transfer for the explicit dim-$6$ loop block]\label{cor:D6_basis_transfer}
Let $B\in GL_2(\mathbb R)$ be any fixed change of basis on the same CP-even dim-$6$ scalar sector and set
\[
D_6^{B}(t,\mu):=B^{-1}D_6^{PRC}(t,\mu)B.
\]
Then
\[
\|D_6^{B}(t,\mu)^{-1}\|_{\mathrm{op}}
\le
\|B^{-1}\|_{\mathrm{op}}\|B\|_{\mathrm{op}}\,C_{6,\mathrm{inv}}
\Big(\frac{u(\mu)}{u(\mu_t)}\Big)^{p_6},
\]
so every fixed basis transported from the concrete plaquette/rectangle/chair block inherits the same polylogarithmic inverse control.
\end{corollary}

\begin{proof}
This is immediate from
\[
(D_6^{B}(t,\mu))^{-1}=B^{-1}(D_6^{PRC}(t,\mu))^{-1}B
\]
and Theorem~\ref{thm:D6_block}.
\end{proof}

\begin{corollary}[Diagonal-block inverse control beyond the scalar $1\times 1$ list]\label{cor:Lambda_extended}
The required small-flow-time invertibility criterion is now available in the following additional diagonal blocks:
\begin{enumerate}[label=(\alph*),leftmargin=2.3em]
\item the traceless-symmetric dim-4 tensor block, with inverse control $\Lambda_T(t)=\|C_T(t,\mu)^{-1}\|_{\mathrm{op}}$ from
Theorem~\ref{thm:cT_block};
\item the concrete plaquette/rectangle/chair CP-even dim-$6$ scalar loop block, with inverse control
\[
\Lambda_6(t)=\|D_6^{PRC}(t,\mu)^{-1}\|_{\mathrm{op}}
\]
from Theorem~\ref{thm:D6_block}, and likewise any fixed basis transported from it by Corollary~\ref{cor:D6_basis_transfer};
\item any finite direct sum of the scalar $1\times 1$ sectors and the new finite blocks, after applying the lower-triangular assembly result below.
\end{enumerate}
In each case the inverse growth is at most polylogarithmic in $t^{-1}$.
\end{corollary}

\begin{proof}
Parts (a) and (b) are exactly Theorems~\ref{thm:cT_block} and \ref{thm:D6_block}, together with Corollary~\ref{cor:D6_basis_transfer}.
Part (c) follows by combining those theorems with Proposition~\ref{prop:LaneA_triangular_assembly} below.
\end{proof}
\paragraph{From diagonal blocks to finite lower-triangular families.}
The graded triangular structure expected from power counting is already built into the abstract local-field construction
(Theorem~\ref{thm:field_algebra_truncation}).
For Lane~A one therefore needs not only isolated diagonal-block nondegeneracy, but also a finite-dimensional assembly statement
showing that once the diagonal blocks are controlled, the whole finite lower-triangular system is controlled as well.

\begin{proposition}[Finite lower-triangular assembly of inverse control]\label{prop:LaneA_triangular_assembly}
Let a finite family of flowed probes / renormalized local operators be ordered so that the corresponding coefficient matrix
$\mathbf C(t,\mu)$ is block lower triangular:
\[
\mathbf C(t,\mu)=\bigl(C_{ab}(t,\mu)\bigr)_{1\le a,b\le r},
\qquad
C_{ab}(t,\mu)=0 \ \text{ for } a<b,
\]
with $r<\infty$.
Write
\[
\mathbf D(t,\mu):=\mathrm{diag}\bigl(C_{11}(t,\mu),\dots,C_{rr}(t,\mu)\bigr),
\qquad
\mathbf L(t,\mu):=\mathbf C(t,\mu)-\mathbf D(t,\mu).
\]
Assume:
\begin{enumerate}[label=(\roman*),leftmargin=2.3em]
\item each diagonal block $C_{aa}(t,\mu)$ is invertible and obeys
\[
\|C_{aa}(t,\mu)^{-1}\|\le \Lambda_a(t,\mu);
\]
\item each off-diagonal block is uniformly bounded for $0<t\le t_0$,
\[
\|C_{ab}(t,\mu)\|\le M_{ab}\qquad (b<a);
\]
\item the number of blocks and their sizes are fixed.
\end{enumerate}
Set
\[
\Lambda_{\max}(t,\mu):=\max_{1\le a\le r}\Lambda_a(t,\mu),
\qquad
M_{\mathrm{off}}:=\max_{b<a} M_{ab}.
\]
Then $\mathbf C(t,\mu)$ is invertible for $0<t\le t_0$ and there exists a constant $C_{\mathrm{tri}}<\infty$,
depending only on the block sizes and the finitely many $M_{ab}$, such that
\begin{equation}\label{eq:LaneA_triangular_inverse_bound}
\|\mathbf C(t,\mu)^{-1}\|
\le
C_{\mathrm{tri}}\,(1+\Lambda_{\max}(t,\mu))^{r-1}\,\Lambda_{\max}(t,\mu).
\end{equation}
By finite-dimensional norm equivalence, the same bound holds (after adjusting $C_{\mathrm{tri}}$) for the max-entry norm
$\max_{i,j}|(\mathbf C(t,\mu)^{-1})_{ij}|$ used in Theorem~\ref{thm:LaneA_closure}.
\end{proposition}

\begin{proof}
Factor
\[
\mathbf C(t,\mu)=\mathbf D(t,\mu)\,\bigl(I+\mathbf N(t,\mu)\bigr),
\qquad
\mathbf N(t,\mu):=\mathbf D(t,\mu)^{-1}\mathbf L(t,\mu).
\]
Because $\mathbf N(t,\mu)$ is strictly block lower triangular with $r$ block rows, it is nilpotent of order at most $r$:
\[
\mathbf N(t,\mu)^r=0.
\]
Hence
\[
\bigl(I+\mathbf N(t,\mu)\bigr)^{-1}
=
\sum_{k=0}^{r-1}(-\mathbf N(t,\mu))^k.
\]
Moreover,
\[
\|\mathbf D(t,\mu)^{-1}\|
\le C_{\mathrm{diag}}\,\Lambda_{\max}(t,\mu)
\]
for a constant $C_{\mathrm{diag}}$ depending only on the fixed block sizes and the chosen matrix norm.
Since each off-diagonal block of $\mathbf L$ is uniformly bounded and the block sizes are fixed,
there exists $C_{\mathrm{off}}<\infty$ such that
\[
\|\mathbf N(t,\mu)\|\le C_{\mathrm{off}}\,\Lambda_{\max}(t,\mu).
\]
Therefore
\[
\|\mathbf C(t,\mu)^{-1}\|
\le
\|\mathbf D(t,\mu)^{-1}\|\,
\sum_{k=0}^{r-1}\|\mathbf N(t,\mu)\|^k
\le
C_{\mathrm{diag}}\,\Lambda_{\max}(t,\mu)\sum_{k=0}^{r-1}\bigl(C_{\mathrm{off}}\Lambda_{\max}(t,\mu)\bigr)^k,
\]
which is bounded by the right-hand side of \eqref{eq:LaneA_triangular_inverse_bound} after absorbing the finite geometric sum into
$C_{\mathrm{tri}}(1+\Lambda_{\max})^{r-1}$.
\end{proof}

\begin{corollary}[Finite mixed-sector closure from diagonal-block control]\label{cor:LaneA_mixed_sector}
Consider any finite family of renormalized local operators whose small-flow-time coefficient matrix is block lower triangular
(as predicted by power counting / Theorem~\ref{thm:field_algebra_truncation}).
If each diagonal block admits a quantitative inverse bound $\Lambda_a(t,\mu)$ and each off-diagonal block is uniformly bounded,
then the full family satisfies the Lane~A inverse-control hypothesis \eqref{eq:LaneA_quant_inverse_bound} with
\[
\Lambda_{\mathrm{inv}}^{\mathrm{full}}(t,\mu)
:=
C_{\mathrm{tri}}\,(1+\Lambda_{\max}(t,\mu))^{r-1}\,\Lambda_{\max}(t,\mu).
\]
In particular, whenever the diagonal-block controls are polylogarithmic in $t^{-1}$, the assembled control
$\Lambda_{\mathrm{inv}}^{\mathrm{full}}(t,\mu)$ is again polylogarithmic.
Thus Lane~A applies not only sectorwise, but also to any finite direct sum / lower-triangular assembly built from the presently stabilized
CP-even scalar, CP-odd scalar, small-loop, and traceless-tensor diagonal blocks, together with the concrete plaquette/rectangle/chair dim-$6$ block and any fixed basis transported from it,
and more generally to any finite truncation once the remaining diagonal blocks are added.
\end{corollary}

\begin{proof}
Apply Proposition~\ref{prop:LaneA_triangular_assembly}; a fixed polynomial in polylogarithms is again polylogarithmic.
\end{proof}

\begin{remark}[Role of the explicit sample certificates]\label{rem:LaneA_sample_certificates}
The explicit CP-even scalar, CP-odd scalar, small-loop, traceless-tensor, and concrete dim-$6$ loop blocks now serve two referee-facing purposes.
First, they audit the one-step multiplier and UV-matching machinery in sectors where everything can be written out concretely.
Second, they provide normalization cross-checks for the canonical block protocol used later.
They are \emph{not} a second breadth theorem parallel to the canonical finite-truncation result below;
the load-bearing widening step is the passage from arbitrary fixed symmetry blocks to arbitrary finite engineering caps.
\end{remark}

\paragraph{Canonical flow-lifted blocks and arbitrary finite truncations.}
The explicit scalar/tensor/loop certificates above are useful audited examples, but the Lane~A inverse problem is not inherently restricted to those examples.
For the graded basis of local operators from Definition~\ref{def:local_operator_spaces}, one chooses a \emph{canonical flowed lift} block-by-block:
replace each curvature factor $F_{\mu\nu}$ by the flowed curvature $G_{\mu\nu}(t)$, each covariant derivative by its flowed counterpart,
project to the same $O(4)\times CP$ symmetry type, and impose the same scale-independent normalization moments that fix the graded triangular renormalization matrix of Definition~\ref{def:graded_Z}.
The point is not merely that the resulting matrix ``looks like the identity'' in some coordinates.
Appendix~\ref{app:normalization_crosscheck} proves a basis-free theorem on the exact-dimension quotient block itself:
the small-flow substitution acts as the identity on that quotient at tree level, and Theorem~\ref{thm:LaneA_block_coefficient_theorem} converts that quotient statement into an explicit coefficient-extraction theorem on every arbitrary exact-dimension block.
The proposition below upgrades that theorem-level tree statement to the quantitative $I+O(u(\mu_t))$ UV theorem used by the generic finite-dimensional diagonal-block criterion.

\begin{proposition}[Basis-free exact-dimension UV theorem for an arbitrary fixed symmetry block]\label{prop:canonical_block_uv_identity}
Fix a symmetry block $\mathfrak s$ in the graded basis of Definition~\ref{def:local_operator_spaces}, with engineering dimension $\Delta(\mathfrak s)$,
finite block size $m_{\mathfrak s}$, and quotient block $\mathscr V_{\mathfrak s}$ equipped with the canonical block norm of Definition~\ref{def:LaneA_canonical_block_norm}.
Let
\[
\Sigma_{\mathfrak s}^{\mathrm{UV}}(t):\mathscr V_{\mathfrak s}\longrightarrow \mathscr V_{\mathfrak s}
\]
denote the exact-dimension diagonal matching operator from the canonical flowed quotient block to the coherently renormalized quotient block at the matching scale $\mu_t\simeq 1/\sqrt t$.
Then
\begin{equation}\label{eq:canonical_block_uv_identity}
\Sigma_{\mathfrak s}^{\mathrm{UV}}(t)=I_{\mathfrak s}+E_{\mathfrak s}(t),
\qquad
\|E_{\mathfrak s}(t)\|_{\mathrm{op,can}}\le C_{\mathfrak s,\mathrm{UV}}\,u(\mu_t)
\end{equation}
for all sufficiently small $t$.
In particular, after shrinking $t_0$ so that $C_{\mathfrak s,\mathrm{UV}}u(\mu_t)\le \tfrac12$ on $(0,t_0]$, the UV block is invertible and
\begin{equation}\label{eq:canonical_block_uv_inv}
\|\big(\Sigma_{\mathfrak s}^{\mathrm{UV}}(t)\big)^{-1}\|_{\mathrm{op,can}}\le 2.
\end{equation}
If one writes $\Sigma_{\mathfrak s}^{\mathrm{UV}}(t)$ in the fixed canonical basis
\[
\mathbf Y_{\mathfrak s,t}(f)
=
\bigl(Y^{(1)}_{\mathfrak s,t}(f),\dots,Y^{(m_{\mathfrak s})}_{\mathfrak s,t}(f)\bigr)^{\mathsf T},
\qquad
\mathbf O_{\mathfrak s}^{\mathrm{ren}}(f;\mu)
=
\bigl(O^{(1),\mathrm{ren}}_{\mathfrak s}(f;\mu),\dots,O^{(m_{\mathfrak s}),\mathrm{ren}}_{\mathfrak s}(f;\mu)\bigr)^{\mathsf T},
\]
then the small-flow-time expansion has the block form
\begin{equation}\label{eq:canonical_block_uv_form}
\mathbf Y_{\mathfrak s,t}(f)
=
C_{\mathfrak s,\mathrm{UV}}(t)\,\mathbf O_{\mathfrak s}^{\mathrm{ren}}(f;\mu_t)
+\sum_{\mathfrak s'\prec \mathfrak s} B_{\mathfrak s\mathfrak s'}(t)\,\mathbf O_{\mathfrak s'}^{\mathrm{ren}}(f;\mu_t)
+\mathbf R_{\mathfrak s,t}(f),
\end{equation}
where $C_{\mathfrak s,\mathrm{UV}}(t)$ is exactly the matrix of $\Sigma_{\mathfrak s}^{\mathrm{UV}}(t)$ in that basis.
\end{proposition}

\begin{proof}
Theorem~\ref{thm:LaneA_block_coefficient_theorem} already identifies the exact-dimension diagonal tree-level matching matrix on
$\mathscr V_{\mathfrak s}$ as the identity in the canonical coordinates, with representative-independence and cap-independence built into the quotient-block formulation.
Thus the exact-dimension diagonal matching operator has tree-level part $I_{\mathfrak s}$.

To upgrade from tree level to the physical AF matching scale, use the insertion decomposition at matching scale (Theorem~\ref{thm:ins_decomp}) together with
Proposition~\ref{prop:Eins_bound}: on a fixed finite block, every exact-dimension diagonal coefficient is analytic in the AF coupling and differs from its tree-level value by $O(u(\mu_t))$.
Because $m_{\mathfrak s}<\infty$, the canonical block norm is equivalent to any coefficient norm on that block, so these coefficientwise $O(u(\mu_t))$ bounds assemble into the operator bound in \eqref{eq:canonical_block_uv_identity}.
The inverse estimate \eqref{eq:canonical_block_uv_inv} is the Neumann-series bound for $I_{\mathfrak s}+E_{\mathfrak s}(t)$.
Finally, once the fixed canonical basis is chosen, the operator $\Sigma_{\mathfrak s}^{\mathrm{UV}}(t)$ is represented by the matrix $C_{\mathfrak s,\mathrm{UV}}(t)$, and the full SFTE takes the block form \eqref{eq:canonical_block_uv_form} with the lower-block terms and remainder supplied by Theorem~\ref{thm:ins_decomp}.
\end{proof}

\begin{corollary}[Canonical block matrix theorem]\label{cor:canonical_block_matrix_theorem}
Under the hypotheses of Proposition~\ref{prop:canonical_block_uv_identity},
the exact-dimension diagonal UV matrix in the canonical basis satisfies
\[
C_{\mathfrak s,\mathrm{UV}}(t)=I_{m_{\mathfrak s}}+E^{\mathrm{mat}}_{\mathfrak s}(t),
\qquad
\|E^{\mathrm{mat}}_{\mathfrak s}(t)\|_{\mathrm{op}}\le C_{\mathfrak s,\mathrm{UV}}\,u(\mu_t),
\]
and therefore obeys the coordinate version of \eqref{eq:canonical_block_uv_inv}.
In particular, the later finite-truncation inverse theorem may be read either intrinsically on the quotient blocks or, equivalently, in the chosen canonical matrices.
\end{corollary}

\begin{proof}
This is just Proposition~\ref{prop:canonical_block_uv_identity} written in the fixed canonical basis.
\end{proof}

\begin{remark}[Scope of the canonical UV identity]\label{rem:canonical_block_uv_scope}
Proposition~\ref{prop:canonical_block_uv_identity} concerns only the \emph{exact-dimension diagonal block}.
It does not assert that lower-dimension mixings vanish.
Those lower-dimension contributions are exactly the strictly lower blocks in \eqref{eq:canonical_block_uv_form};
they are handled later by the lower-triangular assembly theorem.
The point of the proposition is more specific:
once both sides are written in the same complete-contraction basis and the coherent normalization moments are fixed,
the map on the exact-dimension quotient is canonically the identity at tree level.
That is the bookkeeping fact which lets the later truncation theorem be stated without enumerating sectors one by one.
\end{remark}

\begin{theorem}[Canonical symmetry-block inverse control]\label{thm:canonical_block_polylog}
Fix any block $\mathfrak s$ in the graded local operator basis of Definition~\ref{def:local_operator_spaces} and let
$C_{\mathfrak s}(t,\mu)$ be the corresponding coefficient matrix from the canonical flowed lifts to the renormalized local basis at scale $\mu$.
Assume the AF-ball hypothesis along the matching trajectory.
Then for all sufficiently small $t$ the block is invertible and obeys
\begin{equation}\label{eq:canonical_block_polylog}
\|C_{\mathfrak s}(t,\mu)^{-1}\|_{\mathrm{op}}
\le
C_{\mathfrak s,\mathrm{inv}}
\Big(\frac{u(\mu)}{u(\mu_t)}\Big)^{p_{\mathfrak s}},
\qquad
p_{\mathfrak s}:=\frac{4C_{M,\mathfrak s}}{A_-},
\end{equation}
so the inverse growth is again at most polylogarithmic in $t^{-1}$.
\end{theorem}

\begin{proof}
Combine Proposition~\ref{prop:block_mult_bound} with Proposition~\ref{prop:canonical_block_uv_identity} and apply Theorem~\ref{thm:block_polylog} with $C_{\mathfrak s}^{(0)}=I_{m_{\mathfrak s}}$.
\end{proof}

\begin{theorem}[Basis-independent finite-truncation inverse control for the full sharp-local basis]\label{thm:finite_truncation_inverse_control}
Fix $\Delta_{\max}\ge 0$ and order the canonical symmetry blocks occurring in the graded basis of Definition~\ref{def:local_operator_spaces}
with engineering dimension $\le \Delta_{\max}$ by increasing engineering dimension and then by a fixed symmetry-type order.
Let
\[
C^{(\le \Delta_{\max})}(t,\mu)
\]
be the full coefficient matrix from the canonical flowed lifts of that truncation to the corresponding renormalized local basis.
Then, for all sufficiently small $t$:
\begin{enumerate}[label=(\roman*),leftmargin=2.3em]
\item $C^{(\le \Delta_{\max})}(t,\mu)$ is block lower triangular in the chosen grading order;
\item each diagonal block satisfies the polylogarithmic inverse bound of Theorem~\ref{thm:canonical_block_polylog};
\item each off-diagonal block is uniformly bounded on $(0,t_0]$ by the insertion integration and insertion-to-OS bounds (Proposition~\ref{prop:Eins_bound} and Theorem~\ref{thm:ins_to_OS}) together with finite-depth composition in the fixed truncation.
\end{enumerate}
Consequently there exist constants $C_{\Delta_{\max},\mathrm{inv}}<\infty$ and $p_{\Delta_{\max}}<\infty$ such that
\begin{equation}\label{eq:finite_truncation_polylog}
\big\|\big(C^{(\le \Delta_{\max})}(t,\mu)\big)^{-1}\big\|_{\mathrm{op}}
\le
C_{\Delta_{\max},\mathrm{inv}}
\Big(\frac{u(\mu)}{u(\mu_t)}\Big)^{p_{\Delta_{\max}}}
\qquad (0<t\le t_0).
\end{equation}
In particular, the Lane~A quantitative inverse-control hypothesis \eqref{eq:LaneA_quant_inverse_bound} holds for every finite engineering-dimension truncation of the sharp-local operator basis.
\end{theorem}

\begin{proof}
Power counting gives the block lower-triangular structure in the grading order.
Each diagonal block obeys canonical polylogarithmic inverse control by Theorem~\ref{thm:canonical_block_polylog}.
For the finitely many off-diagonal blocks, Proposition~\ref{prop:Eins_bound} and Theorem~\ref{thm:ins_to_OS} give a scale-uniform control on the corresponding insertion coefficients; because the truncation is finite, repeated composition across the finite RG depth only changes the constant.
Apply Proposition~\ref{prop:LaneA_triangular_assembly} to the resulting block lower-triangular system and absorb the finite polynomial in the diagonal polylog factors into a single exponent $p_{\Delta_{\max}}$ and prefactor $C_{\Delta_{\max},\mathrm{inv}}$.
\end{proof}

\begin{corollary}[Lane~A closure on every finite truncation]\label{cor:LaneA_all_finite_truncations}
For every $\Delta_{\max}$, the finite family of all canonical flowed lifts / renormalized local operators of engineering dimension $\le \Delta_{\max}$ carries the intrinsic capwise data used in Theorem~\ref{thm:LaneA_closure}:
\begin{enumerate}[label=(\roman*),leftmargin=2.3em]
\item the quantitative inverse-control package of Theorem~\ref{thm:finite_truncation_inverse_control};
\item the canonical SFTE / remainder-compatibility package of Corollary~\ref{cor:finite_truncation_SFTE}.
\end{enumerate}
Therefore Theorem~\ref{thm:LaneA_closure} applies canonically and without extra sectorwise hypotheses to every finite engineering-dimension truncation of the sharp-local algebra.
\end{corollary}

\begin{proof}
This is exactly the combination of Theorem~\ref{thm:finite_truncation_inverse_control}, Corollary~\ref{cor:finite_truncation_SFTE}, and the finite-cap extension statement of Theorem~\ref{thm:LaneA_closure}.
\end{proof}

\begin{remark}[Current scope after the finite-truncation upgrade]
The explicit scalar/tensor/loop blocks above remain useful audited examples and normalization anchors.
But the canonical-block results now close Packet~5 at the level actually needed downstream:
every fixed symmetry block in the graded basis of Definition~\ref{def:local_operator_spaces} is controlled by the basis-free associated-graded identity theorem and hence has a canonical UV normalization,
every such block therefore satisfies the polylogarithmic inverse bound of Theorem~\ref{thm:canonical_block_polylog},
and every finite engineering-dimension truncation inherits both quantitative inverse control and SFTE/remainder compatibility by Theorem~\ref{thm:finite_truncation_inverse_control} and Corollary~\ref{cor:finite_truncation_SFTE}.
Table~\ref{tab:LaneA_packet_interface_map} in Appendix~\ref{app:normalization_crosscheck} records the exact role boundary: the present section exports the unconditional finite-truncation input used later by Theorem~\ref{thm:finite_cap_flowed_sector_bridge} and Theorem~\ref{thm:LaneA_closure}, but it does not by itself define the capwise mixed state, the inductive-union sharp-local state, or bounded-base cyclicity.
Appendix~\ref{app:normalization_crosscheck} is therefore the referee-facing dictionary for the chosen canonical bases and normalization conventions, not a second closure proof.
\end{remark}



\subsection{Lane A finite-cap extension theorem and inductive-union upgrade}
On the frozen proof spine, Lane~A has two roles.
First, it extends the already-constructed Lane~C continuum state on the flowed observable algebra $\mathcal A_{\mathrm{flow}}$ to the full adjoint-local sharp-local algebra by the theorem-facing packet chain
\[
\begin{aligned}
\text{basis-free exact-dimension quotient blocks}
&\Longrightarrow \text{one-shell transport + finite-truncation inverse / SFTE package}\\
&\Longrightarrow \text{finite-cap flowed-sector bridge}
\Longrightarrow \text{inductive union}.
\end{aligned}
\]
The scalar/tensor/loop sample sectors remain in the paper only as audited illustrations and normalization cross-checks; they are not theorem-level inputs on this chain.
The finite-cap bridge is the first mixed-correlator closure node on that chain, and the inductive-union passage occurs only after Theorem~\ref{thm:LaneA_closure} and Corollary~\ref{cor:LaneA_full_scope}.
Second, it carries the bounded positive-time state used in Appendix~F into the continuum-time OS semigroup on the reconstructed sharp-local Hilbert space and proves that this same bounded base algebra is still cyclic there.
The theorem below is therefore the intrinsic finite-cap inverse-flow/OS extension mechanism, now fed by unconditional finite-truncation inputs: the insertion-RG file provides the basis-free canonical one-shell increment theorem, Appendix~F supplies the finite-truncation SFTE/remainder package, and the coefficient file proves inverse control on every finite sharp-local truncation.
It does not build a second independent continuum state on flowed observables: Theorem~\ref{thm:C3U} already furnishes that state, and Lane~A extends it capwise before Corollary~\ref{cor:LaneA_full_scope} passes to the full inductive-limit sharp-local algebra.
\begin{quote}
\emph{Lane A constructs the sharp local operator algebra and verifies the OS axioms for the limiting local state.
The reflection-positivity input is taken from the unflowed Wilson state (Section~\ref{sec:patching});
no reflection-positivity property is assumed for the auxiliary flowed observable sector.}
\end{quote}

\paragraph{OS-positive base state.}
Fix a weak-window target value $u_{\mathrm{ren}}\in(0,u_\ast]$ and let
\[
a_n:=\ell_0 2^{-n},\qquad \beta_n:=\beta(a_n;u_{\mathrm{ren}})
\]
be the associated tuned dyadic trajectory.
Theorem~\ref{thm:continuum_tightness} furnishes the bounded positive-time continuum Wilson state $\omega_{\mathrm{bd}}^{(u_{\mathrm{ren}})}$ on $\mathcal A_+$ together with the actual dyadic OS matrix elements needed downstream.
This is the only reflection-positive base state used by Lane~A.

\paragraph{Flowed probes.}
Lane~C also provides continuum Schwinger functions for the flowed probe family $\mathcal A_{\mathrm{flow}}$ (Theorem~\ref{thm:C3U}).
Because gradient flow smears in the Euclidean time direction, reflection positivity does not transfer to the full flowed algebra by a support argument,
and we do not use (or require) reflection positivity at positive flow time.
The role of flow in Lane~A is renormalization and identification: it produces UV-regular probes and a small-flow-time expansion that
isolates the sharp local operators.

\paragraph{Already-constructed Lane~C flowed state.}
For the same tuned dyadic trajectory, Theorem~\ref{thm:C3U} already furnishes a unique continuum Euclidean state
\[
\omega_{\mathrm{flow}}^{(u_{\mathrm{ren}})}:\mathcal A_{\mathrm{flow}}\to\mathbb C.
\]
Lane~A does not construct a second state on flowed observables: it keeps $\omega_{\mathrm{flow}}^{(u_{\mathrm{ren}})}$ fixed on $\mathcal A_{\mathrm{flow}}$ and extends it to sharp-local insertions by inverse-flow limits plus quantitative inverse control.
For brevity below, write $\omega_{\mathrm{flow}}:=\omega_{\mathrm{flow}}^{(u_{\mathrm{ren}})}$.

\begin{lemma}[RP-safe positive-time representatives for flowed probes]\label{lem:RP_safe_flow_rep}
Fix $\varepsilon>0$ and let $f$ be supported in $\{x_0\ge \varepsilon\}$.
For $t\in(0,t_\ast]$ and $R\ge 1$ with $R\sqrt t\le \varepsilon/2$, define a truncation operator $T_{R,t}$ on lattice configurations
by replacing the initial field outside the Euclidean neighborhood $B(\mathrm{supp}(f),R\sqrt t)$ by a fixed reference configuration.
Set
\[
\mathcal O^{(R)}_{i,t}(f):=\mathcal O_{i,t}(f)\circ T_{R,t}.
\]
Then $\mathcal O^{(R)}_{i,t}(f)$ is a bounded cylinder functional depending only on underlying variables supported in $\{x_0>0\}$ and hence lies in the unflowed positive-time algebra $\mathcal A_+$.

Moreover, for any bounded $G\in\mathcal A_+$,
\[
\big|\omega_a(\Theta G\,\mathcal O_{i,t}(f)) - \omega_a(\Theta G\,\mathcal O^{(R)}_{i,t}(f))\big|
\le \|G\|_\infty\,\|\mathcal O_{i,t}(f)-\mathcal O^{(R)}_{i,t}(f)\|_\infty.
\]
By the flow quasi-locality bounds (Theorem~\ref{thm:tail_R}), the truncation error satisfies
\[
\|\mathcal O_{i,t}(f)-\mathcal O^{(R)}_{i,t}(f)\|_\infty\ \le\ C\,t^{-5}e^{-cR^2}
\]
(on the good-flow chart regime; Lemma~\ref{lem:bad_event_irrelevant} upgrades this to an unconditional bound in OS/RP pairings).
Choosing $R=R(t)\asymp\sqrt{\log(1/t)}$ as in Corollary~\ref{cor:R_choice} yields an error $O(t^\eta)$ for any prescribed $\eta>0$.

Consequently, every use of reflection positivity and OS inner products in this submission is applied only to elements of $\mathcal A_+$:
whenever a flowed probe appears inside an OS/RP expression, it is understood to be replaced by a positive-time representative
$\mathcal O^{(R(t))}_{i,t}(f)\in\mathcal A_+$ with a negligible error at the relevant scale.
\end{lemma}

\begin{proof}
The inclusion $B(\mathrm{supp}(f),R\sqrt t)\subset\{x_0>0\}$ follows from $R\sqrt t\le \varepsilon/2$.
The expectation difference bound is the trivial $\|G\|_\infty$ estimate.
The stated truncation error is exactly the $L^\infty$ consequence of Theorem~\ref{thm:tail_R}; the power-saving choice of $R(t)$ is Corollary~\ref{cor:R_choice}.
\end{proof}

\begin{lemma}[Bad-event contribution is negligible in OS/RP pairings]\label{lem:bad_event_irrelevant}
Fix $\varepsilon>0$ and let $f$ be supported in $\{x_0\ge\varepsilon\}$.
Let $t\in(0,t_\ast]$ and $R\ge 1$ satisfy $R\sqrt t\le\varepsilon/2$, and let $\mathcal O^{(R)}_{i,t}(f)$ be the positive-time representative from
Lemma~\ref{lem:RP_safe_flow_rep}.
Let $\Lambda$ be the finite plaquette set in the neighborhood relevant for $f$ and set
\[
G_{\Lambda}(\varepsilon):=\bigcap_{p\in\Lambda}\{\theta(U_p)\le\varepsilon\}
\]
(as in Lemma~\ref{lem:local_tail}).
Then for any bounded $G\in\mathcal A_+$,
\begin{align}\label{eq:bad_event_irrelevant_bound}
\big|\omega_a(\Theta G\,\mathcal O_{i,t}(f)) - \omega_a(\Theta G\,\mathcal O^{(R)}_{i,t}(f))\big|
&\le \|G\|_\infty\Big( C\,t^{-5}e^{-cR^2} + 2\,\|\mathcal O_{i,t}(f)\|_\infty\,\mathbb P_{\omega_a}\big(G_{\Lambda}(\varepsilon)^c\big)\Big).
\end{align}
Moreover, Lemma~\ref{lem:local_tail} gives the explicit bound
\[
\mathbb P_{\omega_a}\big(G_{\Lambda}(\varepsilon)^c\big)
\le |\Lambda|\,C_{\mathrm{tail}}(G)\,\varepsilon^{-2d_G/3}\,\exp\big(-c_{\mathrm{tail}}(G)\,\beta\,\varepsilon^2\big),
\qquad d_G:=\dim(G),
\]
valid for $\varepsilon\le \varepsilon_{\mathrm{tail}}(G)$ as in Lemma~\ref{lem:local_tail}.
In the AF window, $u\asymp \beta^{-1}$ along the tuned trajectory, so the stretched-exponential tail is \emph{superpolynomial} in $u$:
for every $M$ there exists $u_0(M)>0$ such that $u\le u_0(M)$ implies
\begin{equation}\label{eq:bad_event_prob_superpoly}
\mathbb P_{\omega_a}\big(G_{\Lambda}(\varepsilon)^c\big)\le u^{2M}.
\end{equation}
Consequently, choosing $R=R(t)\asymp\sqrt{\log(1/t)}$ yields a power-saving $O(t^{\eta})$ truncation error from the good-event term.
For the bad-event term, fix $t\in(0,t_\ast]$ and $\eta>0$, choose $M$ large, and then (along the tuned trajectory) take the continuum limit $a\downarrow 0$
far enough into the AF window so that the corresponding weak coupling satisfies $u(a)^M\le t^{\eta}$.
This makes the second term in \eqref{eq:bad_event_irrelevant_bound} also $O(t^{\eta})$ at that fixed $t$.
In particular, all RP/OS pairings may be evaluated using the positive-time representatives with an unconditional $O(t^{\eta})$ error,
and the subsequent SFTE argument then sends $t\downarrow 0$.
\end{lemma}

\begin{proof}
Decompose the expectation difference into contributions from $G_{\Lambda}(\varepsilon)$ and its complement.
On $G_{\Lambda}(\varepsilon)$, Theorem~\ref{thm:tail_R} gives the $L^\infty$ truncation bound
$\|\mathcal O_{i,t}(f)-\mathcal O^{(R)}_{i,t}(f)\|_\infty\le C t^{-5}e^{-cR^2}$.
On $G_{\Lambda}(\varepsilon)^c$, use the trivial bound
$|\omega_a(\Theta G\,\mathcal O_{i,t}(f)) - \omega_a(\Theta G\,\mathcal O^{(R)}_{i,t}(f))|
\le 2\|G\|_\infty\,\|\mathcal O_{i,t}(f)\|_\infty\,\mathbb P_{\omega_a}(G_{\Lambda}(\varepsilon)^c)$.
This yields \eqref{eq:bad_event_irrelevant_bound}.
The superpolynomial probability estimate \eqref{eq:bad_event_prob_superpoly} follows from Lemma~\ref{lem:local_tail} and the elementary calculus fact that $e^{-c/u}\le u^{2M}$ for all sufficiently small $u$ (with $c>0$ fixed).
\end{proof}

\paragraph{Small-flow-time expansion and inverse mixing.}
Fix a finite sector-basis of flowed probes $\{\mathcal O_{i,t}(f)\}_{i=1}^{N}$ and a corresponding sector-basis of renormalized local operators
$\{\mathcal O^{\mathrm{ren}}_{j}(f;\mu)\}_{j=1}^{N}$.
The small-flow-time expansion (SFTE) asserts, for $t\downarrow 0$,
\begin{equation}\label{eq:SFTE_matrix_LaneA}
\mathcal O_{i,t}(f)
=\sum_{j=1}^{N} C_{ij}(t;\mu)\,\mathcal O^{\mathrm{ren}}_{j}(f;\mu)
+\mathcal R_{i}(t;f;\mu),
\end{equation}
in the sense of mixed Schwinger functions against arbitrary finite collections of renormalized local insertions.
For the closure theorem below we do \emph{not} use bare invertibility of $C(t;\mu)$ alone.
Instead the abstract inverse-flow mechanism is that for the sector family under consideration there exist $t_0\in(0,t_\ast]$ and a scale-dependent quantitative inverse control function
$\Lambda_{\mathrm{inv}}(t;\mu)\ge 1$ such that
\begin{equation}\label{eq:LaneA_quant_inverse_bound}
\max_{1\le i,j\le N}\big|\big(C(t;\mu)^{-1}\big)_{ij}\big|
\le M_{\mathrm{inv}}(\mu)\,\Lambda_{\mathrm{inv}}(t;\mu)
\qquad (0<t\le t_0),
\end{equation}
and such that the remainders become negligible after multiplication by the same inverse factor in all mixed correlator channels relevant to the target local state:
\begin{equation}\label{eq:LaneA_inverse_remainder_compat}
\Lambda_{\mathrm{inv}}(t;\mu)\,
\Big|
\omega\Big(
\big(\mathcal R_i(t;f;\mu)-\omega(\mathcal R_i(t;f;\mu))\big)
\prod_{\ell=1}^k \mathcal X_\ell
\Big)
\Big|
\xrightarrow[t\downarrow 0]{}0.
\end{equation}
Here $\omega$ is only a placeholder for the eventual sharp-local extension functional in this abstract mechanism.
The special case $\Lambda_{\mathrm{inv}}\equiv 1$ recovers the uniformly bounded inverse situation, but the theorem is designed to accommodate the
polylogarithmic inverse growth established later in Section~\ref{sec:sharp_local}.
For the public Lane~A theorem below, however, no provisional extension functional appears in the hypotheses:
Theorem~\ref{thm:finite_truncation_inverse_control} and Corollary~\ref{cor:finite_truncation_SFTE} supply the capwise inverse bound and the remainder compatibility intrinsically on every finite engineering cap.
In particular, Corollary~\ref{cor:finite_truncation_SFTE} rewrites the abstract condition \eqref{eq:LaneA_inverse_remainder_compat}
as a concrete family of capwise OS/core matrix-element bounds, so the proof of Theorem~\ref{thm:LaneA_closure}
never invokes an extension functional before the finite-cap state has been defined.

Define the (centered) inverse-flow approximants
\begin{equation}\label{eq:inverse_flow_approx}
\mathcal O^{\mathrm{app}}_{j}(t;f;\mu)
:=\sum_{i=1}^{N}\big(C(t;\mu)^{-1}\big)_{ji}\,\big(\mathcal O_{i,t}(f)-\omega_{\mathrm{flow}}(\mathcal O_{i,t}(f))\big),
\end{equation}
so the centering in the inverse-flow approximants is taken with respect to the already-constructed Lane~C flowed state.

\paragraph{Lane A extension theorem.}
The next statement is the intrinsic finite-cap Lane~A output needed downstream.

\paragraph{How to read the full-scope statement.}
There are three closure levels, and separating them makes the widened Lane~A theorem much easier to audit.
\begin{enumerate}[label=(\roman*),leftmargin=2.3em]
\item \emph{Finite audited sectors.} The scalar/tensor/loop examples show concretely how inverse control is verified, but after the quotient-block hardening they are illustrations rather than the source of the general breadth claim.
\item \emph{Finite engineering truncations.} Corollary~\ref{cor:LaneA_all_finite_truncations} is the breadth input needed here:
it says that both the quantitative inverse-control package and the SFTE/remainder-compatibility package hold intrinsically on every finite cap of the canonical sharp-local basis.
Its blockwise input is now theorem-level: Appendix~\ref{app:normalization_crosscheck} proves the coefficient-extraction theorem for arbitrary exact-dimension quotient blocks,
so the canonical diagonal UV problem is a literal finite-dimensional statement before the truncation assembly is invoked.
\item \emph{Inductive union.} Corollary~\ref{cor:LaneA_full_scope} then uses the coherence of the normalization scheme to identify the same renormalized fields inside every larger cap and passes to the full sharp-local algebra.
\end{enumerate}
Thus the theorem below should be read as the finite-cap Lane~A extension mechanism with no extra sectorwise assumptions,
while Corollary~\ref{cor:LaneA_full_scope} is the separate inductive-union upgrade to the full sharp-local algebra.

\begin{theorem}[Lane~A finite-cap extension of the Lane~C flowed state]\label{thm:LaneA_closure}
Fix a weak-window target value $u_{\mathrm{ren}}\in(0,u_\ast]$ and let
\[
a_n:=\ell_0 2^{-n},\qquad \beta_n:=\beta(a_n;u_{\mathrm{ren}})
\]
be the tuned dyadic trajectory.
Let $\omega_{\mathrm{flow}}^{(u_{\mathrm{ren}})}$ be the already-constructed continuum Euclidean state on $\mathcal A_{\mathrm{flow}}$ furnished by Theorem~\ref{thm:C3U} along this same tuning, and let
$\omega_{\mathrm{bd}}^{(u_{\mathrm{ren}})}$ be the bounded positive-time continuum Wilson state on $\mathcal A_+$.
Let
\[
\mathcal K_{u_{\mathrm{ren}}}(F,G;s):=\lim_{n\to\infty}\omega_{a_n}(\Theta F\,\tau_s G)
\qquad(F,G\in\mathcal A_+,\ s\in\mathbb D_{\mathrm{dyad}})
\]
be the dyadic OS matrix-element limits furnished by Theorem~\ref{thm:continuum_tightness}.

Fix an engineering cutoff $\Delta_{\max}\ge 0$ and let
\[
\mathcal A_{\mathrm{flow}}^{(\le \Delta_{\max})}\subset\mathcal A_{\mathrm{flow}}
\]
denote the capwise bounded flowed subalgebra generated by the canonical flowed lifts entering that truncation.
Let $\mathfrak A_{\mathrm{loc}}^{(\le \Delta_{\max})}$ be the $\ast$-algebra generated by $\mathcal A_{\mathrm{flow}}^{(\le \Delta_{\max})}$, $\mathcal A_+$, and all canonical renormalized local generators of engineering dimension $\le \Delta_{\max}$.
Let
\[
\mathbf Y_t^{(\le \Delta_{\max})}(f)
=
C^{(\le \Delta_{\max})}(t;\mu)\,\mathbf O_{\mathrm{ren}}^{(\le \Delta_{\max})}(f;\mu)
+\mathbf R_t^{(\le \Delta_{\max})}(f;\mu)
\]
be the canonical finite-cap small-flow-time expansion furnished by Corollary~\ref{cor:finite_truncation_SFTE}, written in the capwise bases
\[
\begin{aligned}
\mathbf Y_t^{(\le \Delta_{\max})}(f)
&=\bigl(Y_{1,t}^{(\le \Delta_{\max})}(f),\dots,Y_{N,t}^{(\le \Delta_{\max})}(f)\bigr)^{\mathsf T},\\
\mathbf O_{\mathrm{ren}}^{(\le \Delta_{\max})}(f;\mu)
&=\bigl(O^{(1),\mathrm{ren}}_{\le \Delta_{\max}}(f;\mu),\dots,O^{(N),\mathrm{ren}}_{\le \Delta_{\max}}(f;\mu)\bigr)^{\mathsf T}.
\end{aligned}
\]
For each capwise basis index $j$, define the centered inverse-flow approximant
\[
\mathcal O^{\mathrm{app},(\le \Delta_{\max})}_{j}(t;f;\mu)
:=
\sum_i \big((C^{(\le \Delta_{\max})}(t;\mu))^{-1}\big)_{ji}\,
\big(Y^{(\le \Delta_{\max})}_{i,t}(f)-\omega_{\mathrm{flow}}^{(u_{\mathrm{ren}})}(Y^{(\le \Delta_{\max})}_{i,t}(f))\big).
\]
Then Theorem~\ref{thm:finite_cap_flowed_sector_bridge}, together with
Theorem~\ref{thm:finite_truncation_inverse_control} and Corollary~\ref{cor:finite_truncation_SFTE}, defines a unique sharp-local continuum state
\[
\omega_{\le \Delta_{\max}}^{(u_{\mathrm{ren}})}
\]
on $\mathfrak A_{\mathrm{loc}}^{(\le \Delta_{\max})}$ extending the restriction of $\omega_{\mathrm{flow}}^{(u_{\mathrm{ren}})}$ to $\mathcal A_{\mathrm{flow}}^{(\le \Delta_{\max})}$.
Equivalently, once the state has been constructed, the formal centered generators of Theorem~\ref{thm:finite_cap_flowed_sector_bridge} are written as
\[
\widehat O_j^{(\le \Delta_{\max})}(f;\mu)
=
O^{(j),\mathrm{ren}}_{\le \Delta_{\max}}(f;\mu)
-\omega_{\le \Delta_{\max}}^{(u_{\mathrm{ren}})}\big(O^{(j),\mathrm{ren}}_{\le \Delta_{\max}}(f;\mu)\big).
\]
Then:
\begin{enumerate}[label=(\alph*)]
\item (\textbf{Sharp local identification}) For each capwise basis index $j$ and test function $f$, the inverse-flow approximants satisfy
\[
\begin{aligned}
\lim_{t\downarrow 0}\, \omega_{\le \Delta_{\max}}^{(u_{\mathrm{ren}})}\Big(
&\mathcal O^{\mathrm{app},(\le \Delta_{\max})}_{j}(t;f;\mu)
\prod_{\ell=1}^{k} \mathcal X_{\ell}\Big)\\
={}&\omega_{\le \Delta_{\max}}^{(u_{\mathrm{ren}})}\Big(
\big(O^{(j),\mathrm{ren}}_{\le \Delta_{\max}}(f;\mu)
-\omega_{\le \Delta_{\max}}^{(u_{\mathrm{ren}})}(O^{(j),\mathrm{ren}}_{\le \Delta_{\max}}(f;\mu))\big)
\prod_{\ell=1}^{k} \mathcal X_{\ell}\Big).
\end{aligned}
\]
for every finite collection of renormalized local insertions $\{\mathcal X_{\ell}\}$ drawn from the same cap.
\item (\textbf{OS axioms for the finite-cap local state}) The Schwinger functions of these renormalized local fields satisfy the Osterwalder--Schrader axioms.
In particular, OS reflection positivity holds on the capwise local positive-time algebra, and the associated field polynomials are realized on the explicit common invariant core of bounded flowed positive-time classes recorded in Lemma~\ref{lem:finite_cap_common_core}.
\item (\textbf{Compatibility with the Lane~C flowed state and the tuned-dyadic bounded base state})
The restriction of $\omega_{\le \Delta_{\max}}^{(u_{\mathrm{ren}})}$ to the capwise flowed subalgebra equals the already-constructed Lane~C state:
\[
\omega_{\le \Delta_{\max}}^{(u_{\mathrm{ren}})}\!\restriction_{\mathcal A_{\mathrm{flow}}^{(\le \Delta_{\max})}}
=
\omega_{\mathrm{flow}}^{(u_{\mathrm{ren}})}\!\restriction_{\mathcal A_{\mathrm{flow}}^{(\le \Delta_{\max})}}.
\]
On the bounded positive-time algebra, the extension agrees with $\omega_{\mathrm{bd}}^{(u_{\mathrm{ren}})}$:
\[
\omega_{\le \Delta_{\max}}^{(u_{\mathrm{ren}})}\!\restriction_{\mathcal A_+}
=
\omega_{\mathrm{bd}}^{(u_{\mathrm{ren}})}.
\]
Moreover, for every bounded $F,G\in\mathcal A_+$ and every dyadic shift $s\in\mathbb D_{\mathrm{dyad}}$,
\[
\omega_{\le \Delta_{\max}}^{(u_{\mathrm{ren}})}(\Theta F\,\tau_s G)=\mathcal K_{u_{\mathrm{ren}}}(F,G;s).
\]
\end{enumerate}
\end{theorem}

\paragraph{Mixed-channel downstream closure node.}
For downstream Lane~A closure statements, the positive-time mixed-channel Cauchy clause enters only through Theorem~\ref{thm:finite_cap_flowed_sector_bridge}(i), with Corollary~\ref{cor:finite_cap_mixed_family_packaging} supplying the finite-family packaging/covariance node used there.
Corollary~\ref{cor:finite_truncation_SFTE} and Lemma~\ref{lem:finite_cap_mixed_channel_reduction} contain the precursor finite-cap reductions; downstream closure statements cite only Theorem~\ref{thm:finite_cap_flowed_sector_bridge}(i).
In the notation fixed there, $K_W$ denotes the associated support neighborhood, $C_W^{\mathrm{ins}}$ packages the fixed-word insertion dependence, and the finite truncation data record the remaining theorem-facing constants; expanding $C_W^{\mathrm{ins}}$ recovers $\|W\|_{\mathrm{bnd}}:=\prod_r \|B_r\|_\infty$ together with the normalized-word factor from Corollary~\ref{cor:ins_to_OS_bounded_word}.

\begin{proof}
Fix $\Delta_{\max}$ and abbreviate
\[
\mathfrak A_{\mathrm{loc}}^{(\le \Delta_{\max})}=:\mathcal A_{\mathrm{loc}}^{\mathrm{cap}},
\qquad
\omega_{\le \Delta_{\max}}^{(u_{\mathrm{ren}})}=:\omega^{(u_{\mathrm{ren}})}.
\]
Write
\[
Y_{i,t}(f):=Y^{(\le \Delta_{\max})}_{i,t}(f),
\qquad
O_i^{\mathrm{ren}}(f;\mu):=O^{(i),\mathrm{ren}}_{\le \Delta_{\max}}(f;\mu),
\qquad
C(t;\mu):=C^{(\le \Delta_{\max})}(t;\mu).
\]
Let
\[
\mathcal B_{\mathrm{cap}}^{(\le \Delta_{\max})}
:=
\ast\!\bigl(\mathcal A_{\mathrm{flow}}^{(\le \Delta_{\max})},\mathcal A_+\bigr)
\]
be the bounded capwise algebra used in Theorem~\ref{thm:finite_cap_flowed_sector_bridge}.

\paragraph{Intrinsic capwise definition of the extension.}
For the construction it is convenient to work first with the abstract centered capwise algebra generated by
$\mathcal B_{\mathrm{cap}}^{(\le \Delta_{\max})}$ and the formal symbols
$\widehat O_j^{(\le \Delta_{\max})}(f;\mu)$.
After fixing the finite family of mixed bounded words used in the construction, Theorem~\ref{thm:finite_cap_flowed_sector_bridge}, together with Corollary~\ref{cor:finite_cap_mixed_family_packaging}, furnishes a unique positive unital functional on that abstract algebra, characterized by the inverse-flow limits against the corresponding finite-family packaged bounded capwise states.
No universal mixed bounded capwise state is presupposed here; the bounded input is exactly that locally fixed finite-family package.
We denote that functional by $\omega^{(u_{\mathrm{ren}})}$.
Once the state has been fixed, the formal centered symbols are written as
\[
\widehat O_j^{(\le \Delta_{\max})}(f;\mu)
=
O_j^{\mathrm{ren}}(f;\mu)-\omega^{(u_{\mathrm{ren}})}\big(O_j^{\mathrm{ren}}(f;\mu)\big),
\]
which is exactly the notation used in the theorem statement.
Thus existence of the capwise state has been reduced to the theorem-level mixed-correlator construction, and uniqueness will follow from the same characterization on the generating set.

\paragraph{(a) Sharp local identification.}
Fix $j$ and $f$ and let $\{\mathcal X_\ell\}_{\ell=1}^k$ be any finite family of renormalized local insertions from the cap.
Write the product $\prod_{\ell=1}^k \mathcal X_\ell$ as the corresponding abstract centered capwise word.
Theorem~\ref{thm:finite_cap_flowed_sector_bridge}(i)--(ii) defines the capwise mixed-correlator functional precisely so that replacing the slot
$\widehat O_j^{(\le \Delta_{\max})}(f;\mu)$ by the bounded inverse-flow representative
$\mathcal O_j^{\mathrm{app},(\le \Delta_{\max})}(t;f;\mu)$ changes the expectation by $o(1)$ as $t\downarrow 0$, against every fixed capwise word in the remaining slots.
For this mixed-channel step we therefore cite only Theorem~\ref{thm:finite_cap_flowed_sector_bridge}(i), whose proof packages the upstream chain Theorem~\ref{thm:ins_to_OS}, Corollary~\ref{cor:ins_to_OS_bounded_word}, Proposition~\ref{prop:finite_cap_bounded_factor_insertion}, Theorem~\ref{thm:finite_truncation_inverse_control}, and Lemma~\ref{lem:finite_cap_mixed_channel_reduction}, together with Corollary~\ref{cor:finite_cap_mixed_family_packaging} for the finite-family packaging/covariance step; the fixed-word dependence is recorded there through $K_W$, $C_W^{\mathrm{ins}}$, and the finite truncation data.
Because the same theorem gives independence of the relative rates in all local slots, one may use a common parameter $t$ for the whole finite insertion family.
After identifying the centered symbol with
$O_j^{\mathrm{ren}}(f;\mu)-\omega^{(u_{\mathrm{ren}})}(O_j^{\mathrm{ren}}(f;\mu))$, this is exactly the convergence claimed in part~(a).
The same characterization fixes all mixed moments of the generating set and therefore gives uniqueness of the capwise extension.

\paragraph{(b) OS axioms for the finite-cap local state.}
We explain why the Schwinger functions of the renormalized local fields satisfy the OS axioms.

\emph{Euclidean invariance and symmetry.}
The finite-family packaged bounded capwise states entering Theorem~\ref{thm:finite_cap_flowed_sector_bridge} inherit the same Euclidean invariance on $\mathcal A_{\mathrm{flow}}^{(\le \Delta_{\max})}$ from Theorem~\ref{thm:C3U}, while Corollary~\ref{cor:finite_cap_mixed_family_packaging} records the translation covariance used in the mixed-channel step and the tuned Wilson limit supplies the bounded $\mathcal A_+$ identities.
The local Schwinger functions are defined from those finite-family packaged states by pointwise inverse-flow limits, so the same algebraic identities persist for the limiting local correlators.
Permutation symmetry is exact at finite lattice spacing for the bounded representatives and therefore survives the limit as well.

\emph{Reflection positivity.}
Fix $\varepsilon>0$ and consider the capwise local positive-time algebra generated by insertions $O_j^{\mathrm{ren}}(f;\mu)$ with $\mathrm{supp}(f)\subset\{x_0\ge\varepsilon\}$.
Let $F$ be a finite linear combination of monomials in such generators.
Choose $t\in(0,t_\ast]$ and $R\ge 1$ with $R\sqrt t\le \varepsilon/2$.
Define $F_{t,R}$ by replacing each local generator $O_j^{\mathrm{ren}}(f;\mu)$ appearing in $F$ by its inverse-flow approximant and then replacing each flowed probe inside that approximant by the RP-safe positive-time representative from Lemma~\ref{lem:RP_safe_flow_rep}.
By construction, $F_{t,R}\in\mathcal A_+$.
Since $\omega_{\mathrm{bd}}^{(u_{\mathrm{ren}})}$ is reflection positive on $\mathcal A_+$ by Theorem~\ref{thm:continuum_tightness}, one has
\[
\omega_{\mathrm{bd}}^{(u_{\mathrm{ren}})}(\Theta F_{t,R}\cdot F_{t,R})\ge 0.
\]
Now let $t\downarrow 0$ and choose $R=R(t)\asymp\sqrt{\log(1/t)}$ as in Lemma~\ref{lem:RP_safe_flow_rep}.
By Lemma~\ref{lem:RP_safe_flow_rep}, replacing flowed probes by their positive-time representatives changes mixed Schwinger functions by $o(1)$.
By part~(a), the inverse-flow approximants converge in every fixed capwise channel to the corresponding centered local insertions.
Hence
\[
\omega_{\mathrm{bd}}^{(u_{\mathrm{ren}})}(\Theta F_{t,R(t)}\cdot F_{t,R(t)})
\xrightarrow[t\downarrow 0]{}
\omega^{(u_{\mathrm{ren}})}(\Theta F\cdot F),
\]
and the limit is nonnegative.
Thus OS reflection positivity holds on the capwise local positive-time algebra.

\emph{Regularity and operator domain.}
Theorem~\ref{thm:field_algebra_truncation} and Lemma~\ref{lem:finite_cap_common_core} provide the explicit dense core
$\mathcal D^{(\le \Delta_{\max})}_{\mathrm{flow}}$ of bounded flowed positive-time classes on which all finite field polynomials are defined by graph-norm limits of bounded dyadic approximants.
Corollary~\ref{cor:finite_truncation_SFTE} identifies those graph-norm limits with the Lane~A inverse-flow representatives, and
Theorem~\ref{thm:finite_cap_flowed_sector_bridge}(iv) identifies the resulting positive-time OS/core matrix elements with the capwise mixed-correlator state just constructed.
Thus the finite-cap local fields are realized on the explicit common invariant core appearing in the theorem statement.
For regularity, the flowed correlators are bounded by the Lane~C moment/cumulant package together with the kernel bounds in Appendices~B--D; the capwise SFTE remainder bounds then pass those estimates to the renormalized local correlators, with only polylogarithmic inverse factors, which are harmless on a fixed cap.
Hence the present finite-cap Schwinger functions satisfy OS$_0$.

\paragraph{(c) Compatibility with the Lane~C flowed state and the tuned-dyadic bounded base state.}
The restriction to $\mathcal A_{\mathrm{flow}}^{(\le \Delta_{\max})}$ is
$\omega_{\mathrm{flow}}^{(u_{\mathrm{ren}})}\!\restriction_{\mathcal A_{\mathrm{flow}}^{(\le \Delta_{\max})}}$
by Theorem~\ref{thm:finite_cap_flowed_sector_bridge}(ii).
The same theorem also gives
\[
\omega^{(u_{\mathrm{ren}})}\!\restriction_{\mathcal A_+}=\omega_{\mathrm{bd}}^{(u_{\mathrm{ren}})}.
\]
Finally, for bounded $F,G\in\mathcal A_+$ and dyadic $s\in\mathbb D_{\mathrm{dyad}}$, Theorem~\ref{thm:continuum_tightness} already records the limiting matrix element
$\mathcal K_{u_{\mathrm{ren}}}(F,G;s)$ along the tuned dyadic trajectory, so the same identity holds for the capwise extension:
\[
\omega^{(u_{\mathrm{ren}})}(\Theta F\,\tau_s G)=\mathcal K_{u_{\mathrm{ren}}}(F,G;s).
\]
This completes the proof.
\end{proof}

\begin{corollary}[Lane~A preserves the tuned-dyadic bounded base state]\label{cor:LaneA_bounded_state_restriction}
Fix $u_{\mathrm{ren}}\in(0,u_\ast]$ and let $\omega^{(u_{\mathrm{ren}})}$ be the full sharp-local state furnished by Corollary~\ref{cor:LaneA_full_scope}.
Then
\[
\omega^{(u_{\mathrm{ren}})}\!\restriction_{\mathcal A_+}=\omega_{\mathrm{bd}}^{(u_{\mathrm{ren}})}
\]
and, for bounded $F,G\in\mathcal A_+$ and dyadic $s\in\mathbb D_{\mathrm{dyad}}$,
\[
\omega^{(u_{\mathrm{ren}})}(\Theta F\,\tau_s G)=\mathcal K_{u_{\mathrm{ren}}}(F,G;s).
\]
\end{corollary}

\begin{proof}
Part~(c) of Theorem~\ref{thm:LaneA_closure} gives the stated restriction and dyadic matrix-element identity on every finite engineering cap.
Corollary~\ref{cor:LaneA_full_scope} passes those compatible identities to the inductive union.
\end{proof}

\begin{remark}[Phase-II seam note for Corollaries~\ref{cor:LaneA_bounded_state_restriction}--\ref{cor:LaneA_full_scope}]
Corollary~\ref{cor:LaneA_bounded_state_restriction} is the bounded-state compatibility theorem only.
The finite-cap sharp-local OS structure belongs to Theorem~\ref{thm:LaneA_closure}(b), while the actual passage to the full sharp-local inductive union belongs only to Corollary~\ref{cor:LaneA_full_scope}.
\end{remark}


\begin{corollary}[Full sharp-local scope by finite truncations and inductive union]\label{cor:LaneA_full_scope}
For each $\Delta_{\max}\ge 0$, let $\mathfrak A_{\mathrm{loc}}^{(\le \Delta_{\max})}$ be the finite truncation of the graded local operator basis from Definition~\ref{def:local_operator_spaces},
and choose the canonical flowed lifts block-by-block as organized by Theorem~\ref{thm:LaneA_associated_graded_identity} and Proposition~\ref{prop:canonical_block_uv_identity}.
Then Corollary~\ref{cor:LaneA_all_finite_truncations} implies that Lane~A applies unconditionally to every $\mathfrak A_{\mathrm{loc}}^{(\le \Delta_{\max})}$.
The resulting finite-truncation sharp-local states are compatible under the natural inclusions
\[
\mathfrak A_{\mathrm{loc}}^{(\le \Delta_{\max})}\hookrightarrow \mathfrak A_{\mathrm{loc}}^{(\le \Delta_{\max}')}
\qquad (\Delta_{\max}\le \Delta_{\max}'),
\]
and therefore define a unique state on the inductive union
\[
\mathfrak A_{\mathrm{loc}}:=\varinjlim_{\Delta_{\max}\to\infty}\mathfrak A_{\mathrm{loc}}^{(\le \Delta_{\max})}.
\]
This state extends $\omega_{\mathrm{bd}}^{(u_{\mathrm{ren}})}$ on $\mathcal A_+$ and satisfies the OS axioms on the full sharp-local algebra generated by flowed curvature composites.
\end{corollary}

\begin{proof}
Fix $\Delta_{\max}$.
By Corollary~\ref{cor:LaneA_all_finite_truncations}, the whole truncation $\mathfrak A_{\mathrm{loc}}^{(\le \Delta_{\max})}$ already satisfies both the quantitative inverse-control and the SFTE/remainder-compatibility inputs of Theorem~\ref{thm:LaneA_closure}.
Hence Theorem~\ref{thm:LaneA_closure} applies to that truncation.
If $\Delta_{\max}'\ge \Delta_{\max}$, the canonical flowed lifts and normalization conditions on the smaller truncation are the restrictions of those on the larger truncation,
so the resulting Schwinger functions agree on the overlap.
Any polynomial in sharp-local generators involves only finitely many basis elements, hence belongs to some finite truncation; define its expectation there.
Compatibility makes this independent of the chosen truncation, and the OS axioms pass to the inductive union because every OS test polynomial is finite.
The bounded-state restriction statement is inherited from Theorem~\ref{thm:LaneA_closure}(c) on each finite cap and therefore on the inductive union.
\end{proof}


\begin{theorem}[Bounded positive-time base algebra is cyclic in the full sharp-local reconstruction]\label{thm:LaneA_bounded_base_cyclic}
Fix $u_{\mathrm{ren}}\in(0,u_\ast]$ and let $\omega^{(u_{\mathrm{ren}})}$ be the full sharp-local state furnished by Corollary~\ref{cor:LaneA_full_scope}.
Let $(\mathcal H_{\mathrm{loc}},\Omega_{\mathrm{loc}},H_{\mathrm{loc}})$ be the OS reconstruction of this state from the full sharp-local positive-time algebra.
Let $\mathcal P_{+,\mathrm{sharp}}$ denote the $\ast$-algebra of finite positive-time polynomials in the renormalized local generators
$\mathcal O^{\mathrm{ren}}_j(f;\mu)$ with test functions supported in $\{x_0>0\}$.
Then:
\begin{enumerate}[label=(\roman*),leftmargin=2.3em]
\item for every $P\in\mathcal P_{+,\mathrm{sharp}}$ there exists a family $P_t^{\mathrm{bnd}}\in\mathcal A_+$ such that
\[
\|[P_t^{\mathrm{bnd}}]-[P]\|_{\mathcal H_{\mathrm{loc}}}\xrightarrow[t\downarrow 0]{}0;
\]
\item the image of $\mathcal A_+$ is dense in $\mathcal H_{\mathrm{loc}}$, equivalently $\Omega_{\mathrm{loc}}=[1]$ is cyclic for the bounded positive-time base algebra inside the full sharp-local reconstruction;
\item the centered bounded classes satisfy
\[
\overline{\mathrm{span}}\{[F-\omega^{(u_{\mathrm{ren}})}(F)]:F\in\mathcal A_+\}=\Omega_{\mathrm{loc}}^\perp.
\]
\end{enumerate}
\end{theorem}

\begin{proof}
Let $P\in\mathcal P_{+,\mathrm{sharp}}$.
Only finitely many sharp-local generators occur in $P$, so there exists an engineering cap $\Delta_{\max}$ such that
$P\in\mathfrak A_{\mathrm{loc}}^{(\le \Delta_{\max})}$.
By Corollary~\ref{cor:LaneA_all_finite_truncations}, Theorem~\ref{thm:LaneA_closure} applies unconditionally on that cap.
Choose $\varepsilon>0$ so that every test function appearing in $P$ is supported in $\{x_0\ge \varepsilon\}$.
For $t\in(0,t_\ast]$ choose $R(t)\ge 1$ with $R(t)\sqrt t\le \varepsilon/2$ and $R(t)\asymp\sqrt{\log(1/t)}$ as in Lemma~\ref{lem:RP_safe_flow_rep}.

For each sharp-local generator $\mathcal O^{\mathrm{ren}}_j(f;\mu)$ appearing in $P$, let
$\mathcal O^{\mathrm{app}}_j(t;f;\mu)$ be the centered inverse-flow approximant from \eqref{eq:inverse_flow_approx},
and let $\mathcal O^{\mathrm{bnd}}_{j,t}(f;\mu)\in\mathcal A_+$ be obtained by replacing every flowed probe inside
$\mathcal O^{\mathrm{app}}_j(t;f;\mu)$ by its RP-safe representative from Lemma~\ref{lem:RP_safe_flow_rep}.
Now define $P_t^{\mathrm{bnd}}\in\mathcal A_+$ by evaluating the same noncommutative polynomial expression as $P$
in the bounded representatives $\mathcal O^{\mathrm{bnd}}_{j,t}(f;\mu)$.
Because $P$ contains only finitely many generators, $P_t^{\mathrm{bnd}}$ is a bounded positive-time cylinder observable for each fixed $t$.

The proof of Theorem~\ref{thm:LaneA_closure}(b) already shows how to pass from convergence of the individual inverse-flow approximants
to convergence of any fixed polynomial in them: expand the finitely many monomials, apply Theorem~\ref{thm:LaneA_closure}(a)
for the SFTE/inverse-flow part, and use Lemmas~\ref{lem:RP_safe_flow_rep} and~\ref{lem:bad_event_irrelevant} to replace flowed probes by RP-safe representatives with an $o(1)$ error.
Applying that argument to the finitely many monomials appearing in $P$, in $\Theta P\cdot P$, and in the mixed products with $P_t^{\mathrm{bnd}}$, one gets
\[
\omega^{(u_{\mathrm{ren}})}(\Theta P_t^{\mathrm{bnd}}\,P_t^{\mathrm{bnd}})\to \omega^{(u_{\mathrm{ren}})}(\Theta P\,P),
\qquad
\omega^{(u_{\mathrm{ren}})}(\Theta P_t^{\mathrm{bnd}}\,P)\to \omega^{(u_{\mathrm{ren}})}(\Theta P\,P).
\]
Therefore
\begin{align*}
\|[P_t^{\mathrm{bnd}}]-[P]\|_{\mathcal H_{\mathrm{loc}}}^2
&=\omega^{(u_{\mathrm{ren}})}\!\big(\Theta(P_t^{\mathrm{bnd}}-P)(P_t^{\mathrm{bnd}}-P)\big)\\
&=\omega^{(u_{\mathrm{ren}})}(\Theta P_t^{\mathrm{bnd}}\,P_t^{\mathrm{bnd}})
-2\Re\,\omega^{(u_{\mathrm{ren}})}(\Theta P_t^{\mathrm{bnd}}\,P)
+\omega^{(u_{\mathrm{ren}})}(\Theta P\,P)
\xrightarrow[t\downarrow 0]{}0,
\end{align*}
which proves~(i).

For~(ii), by construction of the OS Hilbert space of the full sharp-local state, the span of classes $[P]$ with $P\in\mathcal P_{+,\mathrm{sharp}}$ is dense in $\mathcal H_{\mathrm{loc}}$.
Part~(i) approximates each such class by classes of elements of $\mathcal A_+$, hence the image of $\mathcal A_+$ is dense.
Since $[1]=\Omega_{\mathrm{loc}}$, this is exactly cyclicity of the bounded positive-time base algebra.
Finally, for every bounded $F\in\mathcal A_+$,
\[
[F]=[F-\omega^{(u_{\mathrm{ren}})}(F)]+\omega^{(u_{\mathrm{ren}})}(F)\,\Omega_{\mathrm{loc}}.
\]
Thus the span of centered bounded classes together with $\Omega_{\mathrm{loc}}$ is dense, and orthogonal projection to $\Omega_{\mathrm{loc}}^\perp$ gives~(iii).
\end{proof}


\begin{theorem}[Canonical scalar non-triviality witness]\label{thm:LaneA_nontriviality_witness}
Fix $u_{\mathrm{ren}}\in(0,u_\ast]$ and let $\omega^{(u_{\mathrm{ren}})}$ be the full sharp-local state from Corollary~\ref{cor:LaneA_full_scope}.
Let $\mathcal E_{\mathrm{ren}}$ denote the centered CP-even dimension-$4$ scalar field constructed from the canonical scalar block,
and let $\phi_{\ast,S}$ be the positive-time reference test function of Lemma~\ref{lem:LaneA_scalar_reference_norm}.
Then
\begin{equation}\label{eq:LaneA_nontriviality_norm_one}
\omega^{(u_{\mathrm{ren}})}\big(\Theta\mathcal E_{\mathrm{ren}}(\phi_{\ast,S})\,\mathcal E_{\mathrm{ren}}(\phi_{\ast,S})\big)=1.
\end{equation}
Consequently the OS class $[\mathcal E_{\mathrm{ren}}(\phi_{\ast,S})]$ is nonzero in $\mathcal H_{\mathrm{loc}}$,
$\mathcal E_{\mathrm{ren}}$ is not a multiple of the identity,
and the full adjoint-local sharp-local theory is non-trivial in its vacuum sector.
\end{theorem}

\begin{proof}
Fix any engineering cap containing the CP-even scalar block.
By Lemma~\ref{lem:LaneA_scalar_reference_norm}, the centered scalar component $X^S_j(\phi_{\ast,S})$ of the finite-truncation renormalized multiplet
has OS norm $1$ at every dyadic scale $t_j$.
By Theorem~\ref{thm:field_algebra_truncation}, these dyadic approximants converge in $L^2$ and hence in OS norm to the limiting centered renormalized scalar field
$\mathcal E_{\mathrm{ren}}(\phi_{\ast,S})$ inside that truncation.
Passing to the limit in the norm identity therefore yields \eqref{eq:LaneA_nontriviality_norm_one} on the finite cap.
Corollary~\ref{cor:LaneA_full_scope} and the cap-compatibility statement of Lemma~\ref{lem:LaneA_block_cap_compat}
identify the same normalized scalar field inside every larger cap and hence on the full inductive-union algebra.
Thus \eqref{eq:LaneA_nontriviality_norm_one} holds in the full sharp-local theory.
A multiple of the identity has zero centered OS class, whereas \eqref{eq:LaneA_nontriviality_norm_one} shows that
$[\mathcal E_{\mathrm{ren}}(\phi_{\ast,S})]$ has norm one.
Hence the field is not a multiple of the identity and the vacuum sector is non-trivial.
\end{proof}

\begin{remark}[Why the inductive-limit step is only a compatibility argument]\label{rem:LaneA_inductive_compatibility}
The only genuinely new bookkeeping input in Corollary~\ref{cor:LaneA_full_scope} is the coherent choice of basis and normalization.
At each exact-dimension symmetry block, the local basis and flowed basis are written using the same complete-contraction formulas,
and the normalization moments are fixed once and for all.
When a larger engineering cap is introduced, Lemma~\ref{lem:coherent_Z} ensures that the fields already constructed in smaller caps are literally the same fields inside the larger cap.
Hence the inductive-limit passage is a compatibility argument, not a second renormalization problem.
\end{remark}

\begin{remark}[Lane~A interface status]
The role split for Packets~5--7 is frozen in Table~\ref{tab:LaneA_packet_interface_map} of Appendix~\ref{app:normalization_crosscheck}.
Theorem~\ref{thm:finite_truncation_inverse_control} and Corollary~\ref{cor:finite_truncation_SFTE} provide only the unconditional finite-truncation inverse / remainder data;
Theorem~\ref{thm:finite_cap_flowed_sector_bridge} is the first mixed-correlator closure node;
Theorem~\ref{thm:LaneA_closure} and Corollary~\ref{cor:LaneA_full_scope} perform the capwise extension and then the inductive-union upgrade;
and Theorem~\ref{thm:LaneA_bounded_base_cyclic} is the first theorem that adds bounded-base cyclicity.
No later sentence should be read as re-proving that packet chain in prose.
\end{remark}



\subsection*{Packet 9 interface: axioms, field correspondence, and non-triviality}
Section~\ref{sec:sharp_local} contains the constructive Lane~A inputs: RP-safe representatives, SFTE/inverse-flow identification, exact-dimension quotient control, finite-truncation inverse estimates, inductive-union closure, and bounded-base cyclicity in the full sharp-local reconstruction.
On the public theorem spine, the later permitted Lane~A exports to Packet~9 are Theorem~\ref{thm:LaneA_closure}(b), Corollary~\ref{cor:LaneA_full_scope}, Theorem~\ref{thm:LaneA_temperedness_verification}, Theorem~\ref{thm:LaneA_bounded_base_cyclic}, and Theorem~\ref{thm:LaneA_nontriviality_witness}.
Appendix~\ref{app:axioms_package} is Packet~9 on the public theorem docket and repackages the first four of those Lane~A exports together with the Route~1 vacuum-gap theorem into one theorem-level dossier proving the Euclidean OS axioms, positive-time clustering/vacuum uniqueness, the standard Poincar\'e-covariant Wightman/Haag--Kastler reconstruction, and the field correspondence for the gauge-invariant local differential polynomials represented by the sharp-local algebra; Theorem~\ref{thm:nontriviality_package} then consumes the separate witness export to rule out the zero/identity theory.
For a first-pass Lane~A audit, Appendix~\ref{app:referee_audit}, Extract~B in \S\ref{subsec:hard_module_extracts}, and Stage~4 of Appendix~\ref{app:executive_proof_sketch} isolate exactly which quotient-space, transport, inverse-control, temperedness, cyclicity, and witness theorems are load-bearing.
